\documentclass[12pt]{article}




\usepackage{setspace}
\usepackage{geometry}
\geometry{verbose,lmargin=2cm,rmargin=2cm,bmargin=2cm,tmargin=2cm}
\usepackage[authoryear]{natbib}
\usepackage[normalem]{ulem}
\usepackage{amssymb}
\usepackage{amsmath}
\usepackage{array}
\usepackage{caption}
\usepackage{graphicx}
\usepackage{siunitx}
\usepackage[normalem]{ulem}
\usepackage{colortbl}
\usepackage{multirow}
\usepackage{hhline}
\usepackage{calc}
\usepackage{tabularx}
\usepackage{threeparttable}
\usepackage{wrapfig}
\usepackage{adjustbox}
\usepackage{hyperref}
\usepackage{xcolor}
\usepackage{booktabs}
\usepackage{subcaption}
\usepackage{pdflscape}
\usepackage{newpxtext,newpxmath}
\definecolor{dark-red}{rgb}{0.8,0.15,0.15}
\definecolor{dark-blue}{rgb}{0.15,0.15,0.8}
\definecolor{medium-blue}{rgb}{0,0,0.5}
\hypersetup{
	colorlinks, linkcolor={dark-red},
	citecolor={dark-red}, urlcolor={dark-red}
}
\newcommand{\redbtext}[1]{\textbf{\textcolor[rgb]{1,0,0}{#1}}}
\newcommand{\florian}[1]{\textcolor[rgb]{0,0.5,0.99}{florian: #1}}
\newcommand{\bluetext}[1]{\textcolor[rgb]{0,0,1}{#1}}
\newcommand{\greentext}[1]{\textcolor[rgb]{0.5176471, 0.7647059, 0.1372549}{#1}}




% To fix the disappearing of references when MM compiles - try following?
%    \usepackage{biblatex}
%    \addbibresource{filename.bib}
%    \printbibliography








%% TABLES & FIGURES
% Tables / Figures at end of package
% \usepackage{figcaps}
% \printfigures




\newcommand{\rootdir}{..}  % go down one level
\newcommand{\plots}{\rootdir/\rootdir/Output/plots}
\newcommand{\tables}{\rootdir/\rootdir/Output/tables}




% User-defined table options
\newcommand\fnote[1]{\captionsetup{font=small}\caption*{\emph{Notes}: #1}} % footnote in table




%% MATH 
\DeclareMathOperator*{\argmin}{arg\,min}
\DeclareMathOperator*{\argmax}{arg\,max}




%% ENVIRONMENTS
\newtheorem{definition}{Definition}
\newtheorem{remark}{Remark}
\newtheorem{corollary}{Corollary}




% R modelsummary numerical setup
\usepackage{siunitx}
\newcolumntype{d}{S[
	input-open-uncertainty=,
	input-close-uncertainty=,
	parse-numbers = false,
	table-align-text-pre=false,
	table-align-text-post=false
	]}


%Empty footnote marker
\newcommand\blfootnote[1]{%
	\begingroup
	\renewcommand\thefootnote{}\footnote{#1}%
	\addtocounter{footnote}{-1}%
	\endgroup
}




%Vertically align rows
\newcolumntype{L}[1]{>{\raggedright\arraybackslash} m{#1} }




% Assumption
\newtheorem{assumption}{\normalfont\scshape Assumption}




%% TITLE PAGE
\begin{document}
	\onehalfspacing
	
	
	% \title{Spatial Rents, Garage Location, and Competition in the London Bus Market}
	\title{The Geography of Market Power: Firm Location, Spatial Rents, and Their Social Costs}
	\medskip
	
\author{%
    \Large{Marleen Marra}\thanks{Department of Economics, KU Leuven, 
    Leuven, Belgium.}
    \qquad 
    \Large{Florian Oswald}\thanks{Uni Turin ESOMAS and Collegio Carlo Alberto}
}

\makeatletter
\let\@fnsymbol\@arabic
\makeatother


\footnotetext[3]{We thank Martin Bretschneider, Xinyu Dai and Ian Hsieh 
for excellent research assistance. We also thank participants at the TSE 
and UCL Econometrics of Games, Matching and Networks conference, EARIE, 
UCL, IO$^+$, NBER Transportation Market Design Workshop, as well as 
Giulia Brancaccio, Gabriel Ahlfeldt, Milena Almagro, Estelle Cantillon, 
Pierre-Philippe Combes, Francesco Decarolis, Pierre Dubois, Michele 
Fioretti, Emeric Henry, Jakub Kastl, and Adam Rosen for feedback. We are 
very grateful to Les Stitson from the London Omnibus Traction Society for 
compiling a dataset for this study based on their archives and for 
generously answering numerous questions about the history of garage 
ownership in London. Thanks also to Robert Munster, the owner of 
\href{http://www.londonbusroutes.net}{LondonBusRoutes.net}, for answering 
questions and sharing data on bus route characteristics with us. We can 
be reached via e-mail at \texttt{marleen.marra@kuleuven.be} and 
\texttt{florian.oswald@unito.it}.}
	\date{\today}
	%\JEL{}
	%\Keywords{}
	
	\maketitle
	\begin{abstract}
		
		This paper studies spatial rents in markets where firm location determines 
competitive advantage. Using the London bus market---where operators maintain 
proprietary garage networks and bid for route contracts---we develop a 
structural model linking garage value to procurement profits through dead mile 
costs, local monopoly rents, and economies of density. Exploiting the exclusive 
use of garages, we recast the location problem as one where garages choose 
operators, yielding a tractable discrete choice estimator. Estimates reveal 
meaningful trade-offs between spatial rent components. Counterfactual 
simulations identify an efficiency loss of \textsterling 34M--\textsterling 
51M per year from garage hold-out, and patterns consistent with tacit 
market-sharing among operators.
		
	%	This paper estimates a structural model of the market for public bus transportation services in London, which links the value of a bus garage to expected profits from bus route procurement auctions. The model is used to derive the spatial rents associated with owning a garage with a certain centrality in the network of garages and routes, i.e., the benefits of having low transportation costs, being far removed from competitors, and having multiple garages clustered together. We exploit the unique features of the London bus industry to recover these rents using a standard discrete choice estimator, reflecting the highly complex location choice problem of heterogeneous multi-plant firms competing in an economy with spatial rents. %The magnitude of these spatial rents is based on the identity of the operator owning the garage. 
	%	Even a parsimonious specification is shown to capture remarkably well the observed changes in the garage-operator network since the privatization of this industry in 1994, as constructed from archival data and bus spotter sites. Counterfactual simulations reveal an efficiency loss of \textsterling 34M -- \textsterling 51M per year, of 12.3\%-18.5\% of the total procurement cost of providing public bus transportation in London, resulting from operators holding out the sale of their garages. The degree to which profitable garage transactions are ``blocked" is lower in periods when more firms enter the market, consistent with anti-competitive market-sharing behavior.  \\
		%Counterfactual simulations confirm that collusive behavior is partly responsible for the observed degree of spatial clustering of firms, beyond the estimated economies of density from operating routes from nearby garages.
		%This paper studies  firm conduct in the market for public bus transportation services in London, which is an industry where  and the associated spatial rents, 
		% -- i.e. rents which arise if physical distance limits competitive forces. 
		% A structural model links the value of a bus garage to expected profits from bus route procurement auctions, to derive the spatial rents associated with owning a garage with a certain centrality in the network of garages and routes. 
		
		
		\noindent JEL codes: D44, L41, L92, R41\\
		% D22, %firm behavior: empirical analysis
		% D44, %auctions
		% H57, %procurement
		% L13, %Oligopoly and Othe Imperfect Market,
		% L41, %Monopolozation - Horizontal Anticompetitive practices
		% L92, %Railroads and other surface transportation, 
		% R12, %Size and Spatial Distributions of Regional Economic Activity (urban)
		% R41\\ %Transportation: Demand, Supply, and COngestion - Travel Time- Safety and Accidents - Transportaion Noise
		% {\color{blue} (select a few: I'd say D44, L41, L92, R41 see tex for description)} \\
		Keywords: Firm Location Choice, Procurement Auctions, Public Transport, Urban 
		
		
	\end{abstract}
	\newpage
	
	
%% INTRODUCTION
This paper studies the classical trade-off \citep[e.g.,][]{Hotelling1929,
salop1979monopolistic} between the profit gains that arise when firms locate 
close to demand and the intensified spatial competition that results from it. 
We examine this trade-off in an urban setting with building restrictions 
and an industry that exhibits economies of density when multiple plants 
operate in close proximity. We refer to the combination of these spatial 
drivers of profitability---proximity to demand, distance from rivals, and 
density of own operations---as \emph{spatial rents}. Both economies of 
density and spatial competition generate spillovers across locations that 
quickly render the firm location problem computationally intractable.\footnote{With $F$ firms and $L$ locations, allowing for multiple firms at the same location, the number of possible 
assignments is $F^L$. With magnitudes $F = 10$ and $L = 100$, as in our setting, this would exceed $10^{100}$, 
larger than the number of atoms in the observable universe. Even evaluating 
a single firm's entry decision across all locations requires integrating 
over $2^L$ combinations when spillovers make profits at each location 
depend on the full network configuration.
% keeping in mind here that our setting is different in that it is one operator per garage, and we want each operator have at least one garage. in that case you use the inclusion-exclusion formula for surjective functions from L locations to F firms. For F = L = 2 that is 2, not F^L = 4. the correct formula is apparently
% S(L, F) = Σ_{k=0}^{F} (-1)^k C(F,k) (F-k)^L
% yielding S(100, 10) ≈ 9.997 × 10^99 — essentially 10^100 but about 0.027% smaller. :-)
} As a 
result, the empirical literature provides scant evidence on two important 
questions at the intersection of industrial organization and urban policy: 
How do building restrictions and zoning affect firm location decisions and 
the spatial distribution of market power? And do firms protect local monopoly 
rents by staying out of each other's territories, distorting competition 
and shaping the transport geography of cities? We study these questions in the context of public 
bus transportation services in London and show how to exploit features of this market to overcome the dimensionality problem of multi-firm location entry with spatial rents. 

In this market, private bus operators maintain proprietary bus fleets from their own garages and bid for contracts 
to operate bus routes across Greater London. Operators are less competitive 
in route auctions for routes far from their garages, due to the \emph{dead 
mile} costs of transporting empty vehicles between garage and route. On 
average, supplying 100 minutes of contracted bus service requires 13 minutes 
of dead mile driving time, and this has an economically meaningful effect on 
procurement outcomes and firms' bidding behavior. As a result, the spatial layout and composition  
of each operator's garage network is the primary determinant of its 
competitive position and expected profitability---making this market 
unusually clean for the study of spatial rents.\footnote{The setting also restricts 
the set of choice variables for participating firms: TfL specifies the exact 
route, frequency, vehicle type, and contract duration, so that garage 
location remains the main margin on which operators can differentiate 
themselves.} Moreover, the exclusive use of garages by single firms and the 
scarcity of vacant industrial land in central London---only 2 percent as of 
2021---means the garage stock is nearly fixed over time, which, as we show, 
is key to the tractability of the model.

We develop a \emph{reversed-choice} framework that exploits the exclusive 
use of garages to recast the firm location problem as one where each garage 
chooses a single operator, rather than the substantially more complex problem of operators choosing a \emph{set} of garages. This reversal 
is applicable more broadly to markets where locations are used exclusively 
by single firms (or plots are small enough that each accommodates only one 
occupant) and ownership transitions are observed sequentially in the 
data---conditions that characterize many markets for commercial real estate, 
retail space, and industrial land. Under these conditions, the reversed-choice 
framework delivers three advantages over standard approaches. First, it 
sidesteps the equilibrium multiplicity problems that arise when modeling 
simultaneous location choices across many firms: conditioning on the observed 
network at each ownership transition yields a unique prediction for each 
sequential choice. Second, transaction prices cancel from the 
equilibrium assignment condition because they enter symmetrically on both 
sides of the ownership transfer, so neither observing nor estimating them 
is required---particularly valuable since transaction prices for commercial 
and industrial property are typically proprietary and difficult to obtain. 
Third, the framework accommodates unusually rich spatial spillovers across 
locations without requiring the high-dimensional integration over strategies 
that makes standard multi-player entry games computationally intractable, or moment inequalities set-identifying structural features of interest. 
The garage value function is grounded in a structural procurement auction 
model that links spatial positioning to expected profitability through both 
cost and mark-up channels, while remaining agnostic about competitive 
conduct at the auction stage. The tight structural link between garage value 
and expected procurement revenues from TfL route auctions provides a direct 
monetary interpretation of the structural estimates in the absence of garage transaction price data. 

A parsimonious specification with five network variables fits the data 
remarkably well, achieving a McFadden pseudo-$R^2$ of 0.554 and a mean AUC
of 0.868.\footnote{\emph{Area Under the Curve (AUC)} is a standard measure in $[0,1]$ used in evaluating the predictive fit of discrete choice models, summarising how well the estimated utilities rank chosen alternatives above unchosen ones across all operators and time periods.} The structural estimates reveal that an operator is indifferent 
between having its closest own garage 3.7 minutes nearer by road and having 
one fewer competing operator's garage within a ten-minute drive. Incumbency 
benefits accumulate steadily with tenure, consistent with the gradual 
build-up of garage-specific investments, while garage value is otherwise 
determined by spatial positioning rather than historical patterns of use. 
This suggests that the observed clustering of operators in certain parts of 
London reflects genuine spatial economic forces rather than mere inertia.

To assess robustness to the sequential timing assumption, we estimate a 
bundle choice specification that treats pairs of consecutive ownership 
transitions as a joint decision. This increases the dimensionality of the 
choice problem from 1,841 to 33,513 observations and introduces a different 
implicit structural error variance, so direct comparison of coefficients across 
specifications is not informative. We instead test equality of coefficient 
ratios, which are scale-invariant and directly comparable across logit 
models. A delta method test fails to reject equality of the key trade-off 
between economies of density and local monopoly rents ($p = 0.945$), while 
the dead miles coefficient is the one ratio that differs significantly 
($p = 0.019$). Taken together, the baseline and bundle estimates establish that the trade-off between proximity 
to own garages and isolation from competitor garages is the robust central 
finding of the structural estimation, and that this finding does not depend 
on the specific timing assumptions of the estimation framework.

We apply the structural estimates to three sets of counterfactuals. We first
map the spatial rent surface across all candidate locations in London to 
identify where the private returns to new garage development are largest. 
Rather than a monotonic geographic gradient (center to fringe of city), the surface reveals dispersed 
high-value pockets across the city, reflecting that the 2019 garage network 
is broadly well-positioned relative to the route network, with the northeast 
emerging as the main area of genuinely underprovided coverage. 

We then use the model to assess the private returns to garage capacity 
expansion, estimating the value to operators of assigning 50 additional 
buses to each garage individually. The bus-garage assignment equation 
underlying these simulations recovers the correct operating garage for 
46 percent of routes exactly and 62--70 percent within a two-minute 
dead-mileage tolerance against observed tender data. Regressing estimated 
expansion values on structural covariates reveals that the dominant cost 
driver is the average dead miles of currently assigned routes, while 
incumbency tenure and spatial isolation from competitors --- the local 
monopoly rent channel --- are the main drivers of heterogeneity across 
garages. One additional competitor garage within a 10-minute drive reduces 
the return to capacity expansion by \textsterling 735{,}000, approximately 
4--5 percent of the operator-specific baseline expansion value. This 
quantifies a classic divergence between private and social returns to 
infrastructure investment \citep{laffont1993theory}: the socially desirable 
margin of capacity investment would minimize dead miles, reducing 
procurement costs borne by the taxpayer and the environmental burden of 
empty bus travel in a congested city, whereas private incentives are 
additionally governed by market power considerations that generate no such 
benefits. %Correcting this divergence requires regulatory intervention to 
%internalize the externalities that profit-maximizing operators have no 
%incentive to reduce.

Our empirical approach builds on a literature 
that uses endogenous firm location decisions to learn about spatial rents, market structure, and competition 
\citep[e.g.,][]{Mazzeo2002, Seim2006, Jia2008, Holmes2011, 
ellickson2013estimating, nishida2015estimating, zheng2016spatial, oberfield2023plants}. The 
closest antecedents are \citet{Jia2008} and \citet{nishida2015estimating}, 
who estimate location choice games with both economies of density and 
competitive effects, and \citet{Holmes2011} and \citet{houde2023nexus}, who 
quantify economies of density in distribution networks using a moment 
inequality approach. Our 
setting --- many firms, exclusive-use locations, a congested mega-city, and regulated fixed demand --- permits 
identification of an unusually rich set of spatial rents: economies of density, local monopoly rents from strategic 
competitor positioning, and transportation costs simultaneously with a simple estimation strategy.\footnote{For example, \citet{Holmes2011} abstracts from local monopoly rents and agglomeration benefits. On the other hand, the moment inequality approach delivers bounds on the part of fixed costs that varies across locations, which is not identified in the reversed-choice framework. We discuss this limitation when simulating counterfactual (urban planning) policy scenario's.}  %Our setting avoids this: garage locations are fixed and ``demand" equates to expected lump-sum transfers from TfL, so observed operator location choices identify spatial rents directly.  %The moment inequality approach, by contrast, bounds fixed entry costs that %, without requiring moment inequalities. The key is the 
%\emph{reversed-choice} framework: rather than operators choosing garages, 
%garages choose operators, delivering a tractable multinomial discrete choice 
%estimating equation.%\footnote{Garage 
%transaction prices cancel from the equilibrium assignment condition 
%in the reversed-choice framework since they are, in equilibrium, 
%identical across choice alternatives, so estimation does not require 
%observing them. The tight structural link between garage utility and TfL 
%auction revenues then provides a direct monetary interpretation of the 
%spatial rent estimates --- the profitability threshold is determined by 
%contrasting operator entry against vacating the plot to the market --- 
%without the need to bound fixed costs from counterfactual entry decisions.}
Relatedly, \citet{Cao2026} provide a framework to address potential distance endogeneity in location choice models; this concern does not arise here as both garage plots and ``demand" (expected lump-sum payments from Transport for London) are fixed. 

Our setting connects naturally to an urban economics literature on 
agglomeration and the spatial organization of cities 
\citep[e.g.,][]{combes2015empirics, de2011handbook}. 
\citet{heblich2020making} show that railway infrastructure shaped 
workplace-residence separation and land values in London through 
agglomeration forces; we study the complementary question of how industrial 
land frictions and private firm location choices generate market power 
within the city's existing bus route network. The urban policy 
counterfactuals connect to a growing literature on transport network design 
\citep{fajgelbaum2020optimal, allen2022welfare, balboni2019harm, 
kreindler2024optimal} and heterogeneous returns to infrastructure investment 
\citep{brancaccio2024investment}. We identify where private and social 
returns to garage investment diverge, with a single source of a wedge: 
spatial isolation from competitors generates local monopoly rents that push 
operators away from the cost-minimizing configuration. Through this mechanism, private garage ownership transfers rents from the public (taxpayers) to firms, increasing dead miles and consequently inflating congestion, noise, and exhaust emissions beyond what cost minimization alone would require. This connects to the classic urban 
economics of transport investment \citep{small2024economics} and to the 
Pigouvian tradition \citep{pigou1920economics, laffont1993theory} of 
identifying market failures from uninternalized externalities.

The garage value function is grounded in a structural procurement auction 
model \citep{guerre2000optimal, Cantillon2006, Cantillon2007}. Our analysis of location choice externalities 
thus offers a spatial interpretation of a well-established finding in procurement 
design: when quality has noncontractible dimensions, awarding contracts on 
price alone generates adverse selection and distorted investment incentives 
\citep{laffont1986using, bajari2001incentives, lopomo2023optimal, 
decarolis2014awarding, gagnepain2002incentive}.\footnote{TfL procures bus routes on 
price alone conditional on contractible service attributes, leaving garage 
location --- the primary margin on which operators can extract rents above 
cost minimization --- outside the contractual framework. In \cite{marra-oswald}, we highlight this and document that 
the operational efficiency of the bus-garage assignment can improve by 14\%, looking at dead miles alone. The current paper focuses 
on the strategic forces governing the garage ownership network itself.} In a alternative interpretation of the data (Appendix \ref{A:collusion}), we assess whether the spatial clustering of operators 
reflects collusive market-sharing arrangements in addition to 
profit-maximizing location choices, exploiting the structural link between 
garage location values and expected procurement profits to construct implied 
punishment values for unrealized transactions. This contributes to a 
literature on collusion detection in procurement 
\citep[e.g.,][]{porter1983study, Porter1993, Porter1999, bajari2003deciding, 
kawai2022detecting} --- which typically takes firms' assets and locations 
as given when testing for bid-level collusion --- and connects to a 
literature showing how market structure and firm positioning can facilitate 
tacit coordination \citep[e.g.,][]{miller2017understanding, bourreau2021market, 
berry2001mergers, sullivan2017ice}. The spatial pattern of implied punishment 
values correlates strongly with garage isolation, reinforcing the 
private-public wedge in location incentives: garages in more isolated 
positions generate both higher private returns and deter entry by competitors.\footnote{As the expected negative 
relationship of implied minimum ``punishment" (in the collusive interpretation) with the number of active firms is sensitive to specification, we cannot rule out the competitive interpretation and take no stance on whether market-sharing dynamics underlie part of the observed garage network.}

The remainder of the paper is organized as follows. Section~\ref{S:Industry} 
describes the market and data. Section~\ref{S:Model} presents the structural 
model. Section~\ref{S:Results} presents the estimation results and robustness 
checks. Section~\ref{S:Counterfactuals} presents the urban policy 
counterfactuals. Section~\ref{S:Conclusion} concludes. 

%% INDUSTRY AND DATA
\section{Industry and data description}\label{S:Industry}
London Bus Services Ltd (London Buses) is part of Surface Transport within 
\emph{Transport for London (TfL)}, in charge of delivering the Mayor's 
Transport Strategy. This involves transporting over six million passengers on 
7,700 scheduled buses across 675 different routes stretching over 490 million 
kilometers of road per year. Procuring bus operation services from private 
operators costs TfL on average \textsterling 273m annually over the period 
2003--2019, paid for by UK taxpayers. The market has been organized in this 
way since it was privatized along geographic lines in 1994. Today, six large 
international transportation firms and a handful of smaller operators compete 
for tenders.\footnote{The major operators over our sample period are Abellio, Arriva, Go-Ahead, Metroline, RATP, and Stagecoach.} This section describes 
the market structure and our data, documents the spatial clustering of bus 
garages, and provides reduced-form evidence linking garage network position 
to route auction outcomes.

For our analysis of spatial rents, we constructed a comprehensive dataset of 
the operator and bus garage network in London between 1994--2019. Each operator 
maintains its own proprietary garage network, and the location of these garages 
\emph{vis-\`{a}-vis} the routes they service is crucial: operators are less 
competitive in route auctions for routes far from their garages, due to the 
dead mile costs of transporting empty vehicles between the garage and the route 
endpoints. As a result, the distance of an operator's garage to the bus route 
in question has an economically meaningful effect on procurement outcomes and 
firms' bidding behavior, as we document in Section~\ref{S:TfL} below. 
Throughout, we refer to the firm operating buses from a garage as the garage 
\emph{owner}---regardless of whether the firm holds a freehold or long 
leasehold contract---and we exclude the small number of cases where a garage 
is shared between two operators. Figure~\ref{fig:garages-TS-byoperator} documents changes in garage ownership over time by operator in our sample. 

The London Omnibus Traction Society provided a historic dataset based on their 
archive of bus schedules, matching the majority of routes tendered since 2005 
to the garage from which they were operated, and containing information about 
garage openings and closings. We complement this with garage location 
information from the privately-run \href{https://londonbuses.co.uk/}{London 
Bus Routes website}, bus stop locations from the \href{https://tfl.gov.uk/info-for/open-data-users/our-open-data}{TfL Open Data Initiative}, 
and garage addresses verified using the Google Maps API. We use the actual driving time as our preferred measure of distance to fully account for the spatial nature of the problem, including geographic obstacles such as the River Thames that can only be crossed at certain points. For a comprehensive 
description of the market and the data, see \cite{marra-oswald}.

A final dataset contains 2017 cross-sectional ward-level median house prices 
obtained from the UK HM Land Registry Price Paid Data, aggregated to 679 
Greater London electoral wards by the \href{https://data.london.gov.uk/dataset/average-house-prices-by-borough-ward-msoa-and-lsoa-2jkxd/}{Greater London Authority}. Matching each garage to its ward, this 
provides a local land-value control used to address potential confounds in 
the estimated return to capacity expansion in Section~\ref{SS:UP1}.

% stats in the paragraph from function BusLocations:::sample_description_stats()
The structural estimates are based on \emph{movers}, e.g. changes to the ownership of garages observed over time. The estimation sample is based on garage-operator network changes taking place after 1994. We hus exclude the first year when the market was being privatized and when the garage network was being formed, as the initial privatization assignments are not strategic choices in the model's sense. We start with 112 garages and exclude three garages with incomplete information about the ownership history between 1995-2019, seven garages that were shared at some point in time, and one temporary garage. Finally, 10 of the remaining 101 garages never changed ownership in our data, so the estimation is based on (multiple) ownership changes in the remaining 91 garages.

These 91 garages account for a total of 192 ownership changes, which form the basis of our structural estimation. Table \ref{T:garagechangedates} lists the type of ownership change that took place. In 42 percent of the times that the garage-operator network changes, a garage is acquired from another operator. In 32 percent of the time the network changes because an operator occupies a garage for the first time. In the remaining cases the network changes because a garage is closed down (33 times) or temporarily not used to operate bus routes (17 times). We interpret these latter cases as the garage being bought from or sold to \emph{the market}, as described further in the theoretical framework. Furthermore, the relevant unit of time will be a so-called \emph{changedate}, which we define as a date at which we observe any change in the garage-operator ownership structure. There are 114 changedates in our data from 1995 to 2019, with on average 1.68 ownership changes per changedate. In 66\% of cases there is a single operator changing garage on a given changedate. The vast majority of network changes (76\%) also involve just one garage per time. These statistics are documented in Table \ref{T:garagechangedates} and exploited in the estimation strategy. %The changedate structure 
%matters for estimation: garages that change hands simultaneously are treated 
%as choosing among operators at the same point in time (see Section~\ref{S:Model}).

To best capture spatial rents associated with the geographic location of garages relative to routes and other garages, distances are measured as the exact drive times on the London road network.\footnote{Simple straight-line distance between points would miss a great deal of nonlinearities arising from geography and, ultimately, true transportation costs. We therefore use the \href{http://project-osrm.org/}{Open Source Routing Machine (OSRM)} to compute optimal drive times between all garages and all other garages and all 55,578 bus stops in London. Besides the distance to the endpoints of the routes, we also consider the distance to the closest stopover point, where at least 5 routes intersect. See \cite{marra-oswald} for further details.} We use the average drive time to the two endpoints of a route as our primary dead miles measure (\emph{Dead Miles (Start-Stop Minutes)}). This measure reflects that a route needs to be serviced from both directions, so being near only one endpoint is not necessarily optimal. A garage at the midpoint of a long route, for instance, minimizes total dead miles even if it is not adjacent to either terminus. We consider the sensitivity of our results to other dead mile measures, including the distance to the nearest route endpoint and the average distance to a stopover point where at least 5 routes intersect in Sections \ref{S:TfL} and \ref{S:Robustness}. 


\begin{table}
\centering
\caption{Summary of all garage network changes 1995-2019}\label{T:garagechangedates}
\begin{minipage}{0.5\textwidth}
\centering
\input{\tables/ownership-statistics.tex}
\end{minipage}%
\begin{minipage}{0.5\textwidth}
\centering
\input{\tables/combined-share-changes.tex}
\end{minipage}
\fnote{The left panel accounts all the changes in the garage-operator network observed in the data. The right panel documents the number of operators entering a garage and the number of garages entered by an operator simultaneously, that is, on the same date. The right panel is based on all times when operators opened a new garage (the 61 ``New entry" events) or purchased an existing garage from another operator (the 81 ``Ownership change" events) between 1995-2019.}
\end{table}

\begin{figure}[!tb]
\centering
\includegraphics[width=0.9\textwidth]{\plots/routes-garages-operators.pdf}
\caption{Route network and garage locations by operator group.}
\fnote{Garage-route network at the end of the sample (2019). Points denote locations of garages, line segments indicate bus routes, and colors indicate which operator owns the garage or operates the route.\label{fig:garage-locations}}
\end{figure}

In the remainder of this section, we document new facts about the observed garage-operator network and its relation to bidding in route auctions.

\subsection{Spatial clustering of garages}
Figure \ref{fig:garage-locations} shows a snapshot of our dataset in 2019, 
which links routes to the garage where they were operated from. A key 
observation is the spatial clustering of bus garages. Specifically, the figure 
shows that garages of a certain operator group are concentrated in certain areas 
of London, rather than being uniformly spread out across the city. The 
structural model and counterfactual simulations are designed to understand the 
factors driving this observed spatial clustering. The model is based on changes 
in the garage ownership network over time. Garages change ownership between 1--7 times, and no obvious spatial pattern 
can be detected as to which garages change ownership more frequently (Appendix Figure \ref{fig:ownership-count}). This 
makes systematic selection of peripheral versus central garages into the 
estimation sample unlikely. 

One explanation for operating routes from nearby garages is the specialized 
functions that some garages perform within an operator's network. Driver 
training, maintenance, engineering, logistics, and HR are among the functions 
that are often centralized in the larger garages.\footnote{Table 
\ref{T:garage-syn} in the Appendix lists the specialized functions performed 
by select garages within operators' networks, based on information on their 
websites.} Together with the capacity constraints due to building restrictions 
in the densely populated city, these might generate cost synergies from 
operating routes from nearby garages. For instance, a report of the UK 
Industrial Land Commission documents that in 2021 (close to the end of our sample period), only 2 percent of the 
industrial land in central London is left vacant, constraining expansion opportunities of existing garages.\footnote{See 
\href{https://centreforlondon.org/wp-content/uploads/2022/01/CFL-IndustrialLand-v4-1.pdf}{https://centreforlondon.org/wp-content/uploads/2022/01/CFL-IndustrialLand-v4-1.pdf}.}


\subsection{Link between garage network and revenue from route operation}\label{S:TfL}

\begin{table}[!tb]
\centering
\caption{Route auctions and Garage Networks of winning bidders\label{tab:dm-reg-stats}}
\begin{adjustbox}{max width = \textwidth}
\input{\tables/deadmiles-tender-data-summary.tex}
\end{adjustbox}
\fnote{\emph{Dead Miles} denote a distance measure from garage to a certain point on a bus route. In the context of this table (route auctions), this concerns the distance from the \emph{winner's} garage to the auctioned route. \emph{Start-Stop} is the average of distance to start and stop point of the route, \emph{Stopover} refers to a bus stop with at least 5 intersection bus routes (representing the 95-th percentile of the number of route intersections at bus stops). The count for complete records of dead mile measures is inferior to the number of observed auctions (1907) because we observe the operating garage for only 592 out of 711 tendered routes (i.e 83\% of cases). The distance measures further differ because fewer bus stops can be classified as stop-over than starting or end points, for example. Incomplete route characteristics (Route Length, PVR) are a result of incomplete coverage in our Wikia and LBR data sources. Different combinations of those variables in our models below will imply different effective sample sizes.}
\end{table}


\begin{table}[!tb]
\centering
\caption{Explaining the size of accepted bids for London bus routes\label{tab:dm-reg-DM_startstop_minutes}}
\begin{adjustbox}{max width = 0.9\textwidth}
\input{\tables/DM-regs-combined.tex}
\end{adjustbox}
\fnote{Statistical significance: + p $<$ 0.1, * p $<$ 0.05, ** p $<$ 0.01, *** p $<$ 0.001.  All models employ the drivetime distance in minutes between the winner's garage to a certain spatial characteristic of the auctioned route: Columns (1)-(3) averages distances to start and end points of the bus route, column (4) uses distance to route start point only, column (5) uses distance to stopover points. See table \ref{tab:dm-reg-stats} for variable details.}
\end{table}


We next relate winning bids from the publicly available TfL tender data to our 
network variables to support the claim that the placement of a garage in the network of 
routes and other garages is relevant for route auction outcomes. We observe the 
operating garage for 592 out of 698 routes; the remainder are cases where the 
operating garage was not a mandatory field in the tender documentation and was 
left unreported. To the extent that garages are more likely to be reported when 
they are closer to the route---and thus more strategically salient---this 
selection could bias the reduced-form estimates. The structural estimates of 
Section \ref{S:Model} are not subject to this concern, as they are based on 
garage proximity and capacity measures that are universally observable and 
formalized in the assignment equation \eqref{eq:assignment_buses}. We 
subsequently compute the empty driving time for 1,457 out of 1,907 
route-auctions. We regress the winning bid for a route contract on: i) auction 
and route characteristics, ii) our computed measures of the dead miles driving 
time between the route and the garage where the route is operated from, iii) 
the driving time between the operating garage and the closest garage of the same 
operator (\emph{Closest Own Garage}, to capture agglomeration benefits), and 
iv) the dead miles of the closest competitor garage (to capture local monopoly 
rents).

On average, buses drive 13.28 minutes (or 7.78 km) empty each time they move 
between garage and route, as measured by the \emph{Dead Miles Start-Stop} 
variable (see Table \ref{tab:dm-reg-stats}). An additional minute of dead mile driving time increases the winning bid by 
at least \textsterling 19,000, as shown in columns (1)--(2) of Table 
\ref{tab:dm-reg-DM_startstop_minutes}; the estimate is larger when excluding 
auction-level controls for the Peak Vehicle Requirement and the number of 
bidders, which absorb part of the variation in bids. Adding the strategic 
network variables in column (3) renders the dead miles coefficient insignificant. This pattern is consistent with dead miles affecting both costs---pushing bids 
up---and competitive positioning---potentially compressing mark-ups as operators 
near a route also tend to face more local competition. These offsetting forces 
are difficult to separate in a reduced-form regression on winning bids. The 
structural model of Section \ref{S:Model} identifies the net effect on operator 
profits rather than on winning bids, which sidesteps this issue. 

The results also point to a negative correlation between the accepted bid and 
the proximity of the winning bidder's \emph{Closest Own Garage} to the 
operating garage, consistent with cost synergies from a spatially concentrated 
garage network. The coefficient on this variable should be understood to reflect 
a mix of cost effects and mark-up adjustments from being better positioned than 
competitors. The effect of competitor proximity, by contrast, is more cleanly 
interpretable: since the drive time to a competitor's garage does not directly 
affect the winning operator's costs, this coefficient plausibly reflects the 
competitive effect on equilibrium mark-ups alone. The winning bid is about 
\textsterling 43,000 higher for each additional minute of drive time between 
the operating garage and the nearest competitor garage, consistent with 
operators earning higher mark-ups when shielded from nearby rivals.

The results are robust to alternative dead mile measures, as shown in the final 
two columns of Table \ref{tab:dm-reg-DM_startstop_minutes}. The distance to the 
closest route endpoint (column labelled ``Start") is insignificant in the full specification. The distance 
to the nearest stopover point ---where at least five routes intersect--- (column labelled ``Stopover") enters with 
the expected sign but is statistically weaker  and economically smaller, with 
each additional minute increasing the winning bid by \textsterling 9,000. The other coefficients are also stable across distance specifications. 

%% MODEL
\section{Model}\label{S:Model}
This section develops a static two-stage model that formally links the value 
of a firm's location to the spatial rents it can obtain, in order to estimate 
the determinants of those rents from observed location choices. Time is divided 
into periods, each consisting of two stages. At the beginning of each period, operators observe 
the current garage-operator network and draw idiosyncratic cost shocks for each 
route. In stage 1, changes to the garage-operator network occur. In stage 2, 
operators bid on bus route contracts given the updated network. Firms are 
myopic with respect to future changes in the garage-operator network, but 
account for the persistence induced by multi-year route contracts when bidding. 
The process then repeats in the next period.

\subsection{Stage 2: Bidding for routes given a garage network}\label{SS:bidding}
In the second stage, firms compete for route contracts in a first-price sealed 
bid auction given the existing operator-garage network $\Gamma$, which is 
public information.\footnote{While the first-price auction mechanism simplifies 
the exposition, TfL awards route contracts via combinatorial auctions where 
firms may submit route package bids. Abstracting from package bids is without 
loss of generality. Moreover, \cite{Cantillon2007} find limited empirical 
evidence for cost synergies in operating multiple routes at once. They 
hypothesize that garage capacity constraints may limit the cost savings from 
operating at a larger scale, which we explicitly incorporate in the second 
stage of the model.} The cost for operator $i$ to operate route $r$ from 
garage $j$ can be written as
\begin{align}\label{eq:routecost}
c_{ijr}(\Gamma) = \tau\big(\delta_{ijr}(\Gamma), \nu_{ijr}\big)
\end{align}
where $\delta_{ijr}(\Gamma)$ collects the deterministic cost components: 
route-specific variables such as service length and frequency, and factors 
related to the position of garage $j$ in the network of routes and garages 
$\Gamma$, including dead miles between the garage and the route and, if 
economies of density are present, the distance between garage $j$ and other 
garages owned by operator $i$.\footnote{In practice, operator $i$ operates 
route $r$ from its closest garage; we are more precise about this in Section 
\ref{SS:entry}.} The term $\nu_{ijr}$ captures idiosyncratic unobserved cost 
variation that is private information to the firm, varying independently across 
operators, routes, and garages, with $\tau(\cdot)$ increasing in $\nu_{ijr} 
\sim F_{\nu}$.

Operators participating in the auction bid to maximize expected profit from 
operating the route,
\begin{align}\label{eq:exp_prof}
\pi_{ijr}^{route}(\Gamma) = (b_{ijr} - c_{ijr}(\Gamma))\,P(\text{win}|b_{ijr};\Gamma)
\end{align}
where $P(\text{win}|b_{ijr};\Gamma)$ is the probability of winning the auction 
when submitting bid $b_{ijr}$, given the operator-garage network $\Gamma$. To 
accommodate both competitive and non-competitive explanations for the observed 
spatial clustering of garages, we do not impose that all potential bidders submit competitive bids. 
We only impose the minimal restriction that lower-cost firms achieve higher 
expected procurement profits in equilibrium. This assumption holds for the 
unique Bayes-Nash equilibrium in the competitive case, where the setting 
reduces to a standard independent private values first-price procurement 
auction.\footnote{Under standard regularity conditions, the unique symmetric 
Bayes-Nash equilibrium $\sigma^*(c_{ijr})$ solves \citep{guerre2000optimal}:
\begin{align}\label{eq:eqbid}
\sigma^*(c_{ijr}(\Gamma)) = c_{ijr}(\Gamma) + 
\frac{P(\text{win}|\sigma^*(c_{ijr}(\Gamma));\Gamma)}
{P'(\text{win}|\sigma^*(c_{ijr}(\Gamma));\Gamma)}
\end{align}
where $P'(\text{win}|\sigma^*(c_{ijr});\Gamma)$ denotes the derivative of 
$P(\text{win}|b_{ijr};\Gamma)$ with respect to the bid. The mark-up 
decreases in the number of competitors and is compressed when cost 
distributions are less dispersed, for instance when all operators have garages 
at similar distances from the route.} It also holds for a strong cartel with 
side-payments and, by the results of \cite{Pesendorfer2000study}, for a weak 
cartel without side-payments when sufficiently many contracts are auctioned, 
as both cases yield (approximately) efficient allocations. This assumption 
therefore accommodates a wide range of strategic behavior at the auction stage, 
including bidding cartels with unspecified internal organization.\footnote{For 
example, some firms may form a bidding ring and hold an internal knockout 
auction to determine their representative bidder in the target auction. In that 
case, \eqref{eq:exp_prof} applies to the representative ring member and the 
competitive bidders, and side-payments must be made conditional on winning to 
other ring members \citep[see e.g.,][]{mcafee1992bidding,Pesendorfer2000study,
Asker2010study}.} Where multiple equilibria exist, we assume that players 
coordinate on a symmetric bidding strategy $\sigma^*(\cdot)$, and denote the 
resulting expected equilibrium profits by $\pi^{*route}(\cdot)$.

Economies of density and dead miles affect $\pi^{*route}(\cdot)$ directly through 
costs, but also generate local monopoly rents through the equilibrium mark-up. 
To illustrate this, consider the win probability in the IPV first-price auction 
with additively separable $\tau(\cdot)$:
\begin{align}\label{eq:strucvalue1}
P(\text{win}|b_{ijr};\Gamma) = \prod_{h\neq i} 
\left[1 - F_{\nu}\!\left(\sigma^{-1*}(b_{ijr})-\delta_{hj'r}(\Gamma)\right)\right]
\end{align}
% inside the square bracket is the probability of event "operator h, using his garage j', which implies deterministic costs of \delta(h,j',r), has a cost shock \nu which makes his total cost higher than mine." hence, he submits a higher big, hence I win.
where $j'$ is the garage from which operator $h$ would operate route $r$. 
Conditional on $b_{ijr}$, the probability of winning increases in the distance 
between competitors' garages and the route, since greater distance implies 
higher deterministic costs $\delta_{hj'r}(\Gamma)$ for competitor $h$. Costly 
dead miles therefore generate local monopoly rents, particularly for operators 
whose garages are more isolated from those of their rivals. Symmetrically, if 
agglomeration benefits are present, operating from a denser own garage network 
lowers costs and raises the probability of winning, further increasing expected 
profits. Having established how garage location drives expected rents from TfL 
route auctions, we next turn to what this implies for the garage-operator match 
in the first stage, where $\pi^{*route}(\cdot)$ serves as the value function that 
governs location choices.


 %Buses can be considered stand-alone units with their schedules and performance overseen by the firm's operations centre. %\footnote{For example, if a bus is running late on its own schedule, the operations center controllers contact the driver to take action, as explained on \href{https://www.reddit.com/r/london/comments/15n71jf/bus_drivers_of_london_how_does_the_job_work/}{this forum}: \emph{``Every route has time cards, you leave at a specific time and have multiple way points where you are supposed to be at a predetermined time. London zoo at 9.27, Baker street at 10.22 for example. If you run too late behind schedule then the controllers will instruct you to change the destination board, terminate early and give passengers transfer tickets. You turn the bus at a specific point and then hopefully you’re on time heading the opposite direction again.".}} 

\subsection{Stage 1: Garage Entry}\label{SS:entry}
The value of garage $j$ to operator $i$, $\Pi_{ij}$, is the sum of expected 
profits generated through route procurement across all routes the garage can 
service. Given the role of capacity constraints in the London bus market, we 
must first determine the set routes which are served from each garage. Let $\text{Capacity}_{j}$ 
denote the number of buses that garage $j$ can accommodate, and let 
$\text{PVR}_{r}$ denote the Peak Vehicle Requirement of route $r$---the required 
number of buses needed to operate the route at peak demand, as mandated by TfL. Buses are allocated 
to the garage starting from the nearest routes, until the capacity constraint 
binds. Specifically, letting $r_j^{n:R}$ denote the $n^{th}$ closest out of 
$R$ routes to garage $j$, the sequence of routes ordered by increasing dead 
miles to garage $j$ is
\begin{equation}\label{eq:assignment_buses}
\mathbf{r}_j = \{r_j^{1:R},r_{j}^{2:R},\dots,r_j^{R:R}\}
\end{equation}
with $r_j(m)$ denoting the $m^{th}$ route in this sequence and $PVR_j(m)$ the 
corresponding PVR. We maintain that the profit function exhibits constant 
returns to scale in PVR, so that each bus stationed at a garage contributes 
equally to the value of that garage.\footnote{This assumption is supported by 
the operational structure of the London bus market, where individual buses 
operate on fixed schedules independently monitored by the firm's operations 
centre, with no material scale economies at the vehicle level. For example, if a bus is running late on its own schedule, the operations center controllers contact the driver to take action, as explained on \href{https://www.reddit.com/r/london/comments/15n71jf/bus_drivers_of_london_how_does_the_job_work/}{this forum}: \emph{``Every route has time cards, you leave at a specific time and have multiple way points where you are supposed to be at a predetermined time. London zoo at 9.27, Baker street at 10.22 for example. If you run too late behind schedule then the controllers will instruct you to change the destination board, terminate early and give passengers transfer tickets. You turn the bus at a specific point and then hopefully you’re on time heading the opposite direction again.".}}

Garages are more valuable to operator $i$ when they have greater capacity, are located closer to 
more profitable routes, closer to other garages owned by $i$, and further away from 
competitors' garages. The garage value function that captures these forces is
\begin{align}\label{eq:garagestruc}
\Pi_{ij}(\Gamma) = \sum_{m=1}^{R} \sum_{v=1}^{PVR_j(m)} \beta_{r_j(m)} 
\frac{\pi^{*route}_{ijr_j(m)}(\Gamma)}{PVR_j(m)} \mathbf{1}\!\left(\sum_{l=0}^{m-1}
PVR_j(l)+v \leq \text{Capacity}_j\right)
\end{align}
with $PVR_j(0)=0$, $\mathbf{1}(\cdot)$ the indicator function, and 
$\beta_{r_j(m)}$ the discount rate associated with the timing of the auction 
for the $m^{th}$ closest route. The fraction times $\beta_r$ 
%\frac{\pi^*_{ijr_j(m)}(\Gamma)}{PVR_j(m)}$ 
gives the present value of expected 
profit per bus stationed in garage $j$ from the procurement contract for route $r=r_j(m)$. 
The summation over routes and buses (vehicles $v$), truncated by the indicator function, 
accounts for filling the garage to capacity with buses from routes in order of 
increasing dead miles.\footnote{The garage value function naturally implies 
that operators only exert competitive pressure on routes not too far from their 
garages when the effect of dead miles on operating cost is sufficiently large, 
so that the model does not require an ad hoc distance cutoff in bidding.} 
Distance from competitors' garages generates local monopoly rents through the 
equilibrium bid mark-up, while proximity to routes and own garages affects 
variable profits through operating costs.

Two considerations motivate a semi-structural rather than fully structural 
approach to estimating \eqref{eq:garagestruc}. First, while the route auctions are modeled 
independently for tractability, TfL in practice auctions routes in small 
batches using package bidding. Identifying the joint cost distribution in a 
combinatorial auction model requires observing all bids for each route 
\citep{Cantillon2007}, which is not feasible as the available data contains 
only the winning and lowest bids. Second, taking the garage value 
function directly to data requires taking a stand on competitive conduct at the 
auction stage. Assuming competitive bidding would preclude the assessment of 
whether collusive market-sharing explains the observed spatial clustering, while 
characterizing the auction equilibrium under collusion requires strong 
assumptions about the cartel's internal organization that are not warranted 
without further evidence. 

We therefore follow the standard convention of the entry literature and adopt 
a semi-structural approach \citep[e.g.,][]{Seim2006, Jia2008, Holmes2011, 
ellickson2013estimating}, characterizing $\Pi_{ij}$ as a function known up to 
a vector of parameters but informed by the structural representation in 
\eqref{eq:garagestruc}. Modeling the auction stage explicitly serves two 
purposes even within this semi-structural framework. First, it disciplines the 
choice of network variables by separating those that operate through costs---such 
as dead miles and own garage proximity---from those that operate through 
equilibrium mark-ups---such as competitor garage proximity---as illustrated in 
Section \ref{S:TfL}. Second, it supports the counterfactual welfare calculations 
in Section \ref{S:Counterfactuals}, where we link estimated location utilities 
to changes in procurement costs.

We now describe how the network $\Gamma$ is formed. A salient feature of the 
London bus garage market is that garage use is exclusive: the incumbent operator 
must be willing to vacate a garage before another firm can enter. In addition, building 
constraints in London mean that the set of garages is nearly constant over 
time. These two features---exclusive use and a stable garage stock---allow us 
to recast the entry game as one where \emph{each garage choose one operator}, rather 
than operators choosing a \emph{set} of garages. As we show in the next section, this reversal delivers a tractable multinomial estimation 
framework and avoids the need to solve explicitly for 
operators' location strategies. Moreover, since all potential entrants must pay the market garage price, prices cancel in the garages choose operators framework, and we 
do not need to observe them to estimate the model, as formalized below.

\subsubsection{Garages choose operators}
We assume complete information at this stage. This is motivated by the 
repeated nature of interactions among the same set of firms across all garage 
transactions: operators are well-informed about each other's valuations, and 
buyout decisions involve extended negotiation rather than simultaneous moves. 
Under complete information, the garage price determination process can be 
described as a second-price auction, as in other negotiated-price asset markets 
\citep{allen2019search, Slattery2021}: the winning operator pays a price equal 
to the value of the garage to its second-most-interested buyer, since the 
incumbent will not sell below that threshold and competing operators will not 
bid above their own valuation.

The equilibrium price of garage $j$ is therefore
\begin{equation}\label{eq:eqprice}
p_{j}^* = \Pi_{j}^{N-1:N}
\end{equation}
where $\Pi_{j}^{N-1:N}$ denotes the second-highest value of garage $j$ among 
all $N$ market participants, with $\Pi_{ij}(\Gamma)$ reflecting the expected 
profits to operator $i$ in the hypothetical scenario where $i$ enters garage 
$j$, holding the rest of the operator-garage network fixed. All potential 
entrants face this same price: $p_{ij} = p_{j}^*$. Importantly, the incumbent 
operator is included in this order statistic, since the highest-value entrant 
can only acquire the garage at a price exceeding the incumbent's own valuation. 
The transaction price thus serves as an endogenous opportunity cost of 
retention for the incumbent, analogous to a second-price auction with the 
seller as an additional bidder. Since transaction prices enter the indirect utility functions of all operators symmetrically, they cancel in equilibrium, and we do not 
need to observe them to estimate the model. We further assume that operators 
are driven by their full variable profits in this price negotiation process (including the 
idiosyncratic component $\epsilon_{ij}$), supporting that at least the identity of the 
highest-value operator for each garage is common knowledge. This is consistent 
with the repeated and transparent nature of garage transactions among a small 
set of firms.

To formally characterize the equilibrium of the garage-chooses-operator game, 
the payoff for garage $j$ under action $a_j = i$ is the value to operator $i$ 
net of the transaction price:
\begin{equation}\label{eq:garageutility}
U_{j}(a_{j},\mathbf{a}_{-j}) = \Pi_{ij}(\Gamma) - p_{j}^*(\mathbf{a}_{-j})
\end{equation}
where we write $p_{j}^*(\mathbf{a}_{-j})$ to make explicit that the 
equilibrium transaction price depends on the actions of all other garages 
through their effect on the network $\Gamma$. Since $p_{j}^*(\mathbf{a}_{-j})$ is determined by the second-highest 
valuation among all operators, the winning operator --- whose dominant 
strategy in the second-price auction is to bid their true valuation --- 
takes this price as given and cannot influence it by their own action. 
The transaction price therefore cancels from the best-response 
correspondence, and the equilibrium assignment reduces to 
\begin{equation}\label{eq:garagechoice}
a_{j}^* = \argmax_{i \in \mathcal{N}} \Pi_{ij}(\Gamma)
\end{equation}
so that each garage is assigned to the operator with the highest value 
for it, and transaction prices need not be observed to estimate the 
model.\footnote{The cancellation follows from $p_j^*(\mathbf{a}_{-j})$ 
not varying with $i$ in the maximisation:
\begin{equation}
a_j^* = \argmax_{i \in \mathcal{N}} U_j(i, \mathbf{a}_{-j}) 
= \argmax_{i \in \mathcal{N}} \left[\Pi_{ij}(\Gamma) - p_j^*(\mathbf{a}_{-j})\right] 
= \argmax_{i \in \mathcal{N}} \Pi_{ij}(\Gamma)
\end{equation}
Note that agglomeration benefits make garages potential complements for the same operator, which by \citet[Theorem~4]{Milgrom2000} 
could preclude the existence of a competitive equilibrium price \emph{vector} 
across all garages simultaneously. This concern does not arise here because 
prices are determined bilaterally for each garage via \eqref{eq:eqprice}: 
existence of $p_j^*$ requires only a well-defined second-highest valuation 
for each garage, which holds under the maintained assumptions on $\Pi_{ij}$. 
The same argument applies to bundle transaction prices in Section 
\ref{sec:seq-bundle}.}

The set of garages choosing an operator is $\mathcal{C}$ (with cardinality 
$C$). Each garage chooses an operator from the set $\mathcal{N}$ (with 
cardinality $N$) of active operators, with $a_{j} = i \in \mathcal{N}$ 
indicating the selected operator. The vector $\mathbf{a} \equiv \{a_{j} : j = 
1, \dots, C\}$ collects the actions of all garages, and $\mathbf{a}_{-j}$ 
those of all garages other than $j$. The Nash equilibrium of this game is the 
$C$-tuple $\mathbf{a}^*$ such that for every garage $j$:
\begin{equation}\label{eq:garagechoice}
a_{j}^* = \argmax_{i \in \mathcal{N}} U_{j}(i,\mathbf{a}_{-j}^*)
\end{equation}
\sloppy A pure-strategy Nash equilibrium exists when the garage value function 
$\Pi(\cdot)$ satisfies supermodularity---so that operator $i$ entering 
garage $j$ increases $i$'s value for all other garages that become closer 
to its network---and decreasing differences---so that the entry of an 
additional competitor within proximity of garage $j$ reduces its value to 
operator $i$. Under these conditions, a pure-strategy Nash equilibrium 
exists and can be computed via an iterative best-response algorithm 
\citep{topkis1978minimizing, topkis1979equilibrium}. Multiple equilibria 
generally exist in entry games of this type 
\citep[see][]{aguirregabiria2016empirical}. Following \citet{Jia2008} and 
\citet{nishida2015estimating}, among others, we compare structural estimates 
across equilibria to assess robustness; details are provided in 
the next section. 

%{\color{blue} [ taken out from 2.2 -- can now drop all tilde's! work info in before eq 13 ]
% We specify a functional form that is additive 
%in observables and a residual error term $\epsilon_{ij}$:
%\begin{equation}\label{eq:garagevalue}
%\Pi_{ij}(\Gamma) = f\big(X_{ij},G_{ij};\theta\big) + \epsilon_{ij}
%\end{equation}
%}

\subsection{Estimation and identification}\label{SS:estimation}
 We specify a functional form of the garage utility function that is additive 
in observables and a residual error term:\footnote{The structural estimates are in units of utility; their monetary 
interpretation is deferred to Section \ref{S:Counterfactuals}, where we link 
estimated utility differences to revenues from TfL route procurement contracts.}
\begin{align}\label{eq:garagevalue_detail}
\tilde{\Pi}_{ij}(\Gamma) = &\; \underbrace{\theta_1 \text{YearsInc}_{j} + \theta_2 \text{Nonvacant}_i 
+ \theta_3 \text{log(NrGarages)}_i }_{X_{ij}'\theta} \notag \\
&+ \underbrace{\theta_4 \text{RouteStarts10Min}_{j} + \theta_5 \text{MindurOwnGar}_{ij} 
+ \theta_6 \text{CompGar10Min}_{ij}}_{G_{ij}'\theta} + \epsilon_{ij};
\end{align}
where $\theta$ collects the parameters of interest---principally those governing 
economies of density, local monopoly rents, and incumbency benefits---$X_{ij}$ 
contains observed characteristics of the garage and operator, and $G_{ij}$ 
contains network variables describing: the centrality of garage $j$ in the 
exogenous route network; the centrality of garage $j$ in competitors' garage 
networks if operator $i$ were to enter; and the centrality of garage $j$ in 
operator $i$'s own garage network upon entry. The error term $\epsilon_{ij}$ 
varies randomly across operators and garages. It aggregates idiosyncratic route 
operation cost draws $\nu_{ijr}$ from \eqref{eq:routecost} across all routes 
near garage $j$, as well as measurement error arising from the low-dimensional 
representation ($G$) of the network variables ($\Gamma$). The deterministic part of the garage utility function is denoted 
$f(\cdot)$. Next, we discuss our main specification; alternative specifications are considered in Section \ref{S:Robustness}.

The variable $\textit{YearsInc}_{j}$ captures \emph{persistence} in 
this market arising from the multi-year nature of route contracts. The linear 
specification implies a constant rate of decay in the incumbency advantage 
over time, with $\theta_1$ measuring the magnitude of this advantage per 
additional year of tenure. $\textit{Nonvacant}_i$ is an intercept common to all active operators capturing 
the differential value of the garage plot when used for bus operations relative 
to alternative uses---that is, the average utility relative to \emph{the 
market} as outside option. $\log(\textit{NrGarages})_i$ captures potential scale effects in the total number of 
garages an operator owns, with diminishing returns to garage expansion: each 
additional garage contributes less to operator utility as the network grows 
larger. $\textit{RouteStarts10Min}_{j}$ is the number of route endpoints within a 
10-minute drive from garage $j$, capturing \emph{dead miles}: a higher value 
implies less empty driving time between the garage and nearby routes, regardless 
of who operates them. $\textit{MindurOwnGar}_{ij}$ is the minimum drive time 
between garage $j$ and any other garage owned by operator $i$, capturing 
\emph{economies of density} net of any business-stealing effects. 
$\textit{CompGar10Min}_{ij}$ is the number of competing operators' garages 
within a 10-minute drive from garage $j$, capturing \emph{local monopoly 
rents}. 

Some garages are not yet developed at the start of the sample in 1995, some 
remain temporarily vacant, and a few are permanently closed before 2019. To 
avoid modelling garage development and to reflect the highly constrained London 
property market, the garage set $\mathcal{J}$ is held constant over time and 
contains all garages ever in use during the sample period. When a garage is not 
owned by a bus operator it is assigned to \emph{the market}---the outside 
option $i=0$, always included in $\mathcal{N}$---representing the garage plot 
at its alternative use value. Since the level of $\tilde{\Pi}_{ij}$ is not 
identified, we normalise the mean utility of the outside option to zero and 
interpret coefficients accordingly.

The estimation strategy compares $\tilde{\Pi}_{ij}(\Gamma)$ across operators 
$i$ for the same garage $j$, so that garage-specific variables cancel from 
the choice probabilities and their coefficients are not separately identified. 
%All coefficients are identified from utility differences relative to the 
%outside option $i=0$, which represents the garage plot at its alternative use 
%value to the market. 
We nonetheless include $\textit{RouteStarts10Min}_{j}$ 
and $\textit{YearsInc}_{j}$ in \eqref{eq:garagevalue_detail} by normalizing 
them to zero for the outside option. This normalization rests on two 
identifying assumptions.

First, the alternative use value of a garage plot does not depend on its 
proximity to bus routes: a vacant plot repurposed for logistics or light 
industrial use has no direct need for bus route access. The assumption 
requires that the market values the plot for its general locational 
attributes---size, road and rail access, central location---rather than 
bus proximity specifically. If central plots happen to be close to bus routes 
and the market values them partly for this reason, $\theta_4$ would be 
attenuated. Second, the incumbency advantage from tenure is specific to bus 
operations. Bus-operation-specific investments---tailored maintenance 
facilities, driver logistics, depot infrastructure---depreciate over time 
and are of limited relevance to an alternative use of the plot, so that 
$\textit{YearsInc}_j = 0$ for $i=0$ is appropriate. Under both assumptions, 
$\theta_4$ and $\theta_1$ are identified from a difference-in-differences 
argument: each coefficient measures how much more likely a garage is to 
remain occupied by a bus operator---rather than revert to the market---when 
it has low dead miles or was recently acquired, relative to what general 
locational attributes alone would predict. By contrast, garage capacity and transaction prices are garage-specific constants that 
enter identically for operators and the market alike and are therefore not identified from utility 
differences in the reversed-choice framework.\footnote{One practical implication is that garage 
transaction prices need not be observed or estimated to determine the 
counterfactual network allocation, since they cancel from the equilibrium 
assignment condition as shown in Section~\ref{SS:entry}. The same goes for garage capacity, although this is observed for two-thirds of the garages. The assumption that capacity is orthogonal to the operator-specific network 
variables is verified empirically in Appendix~\ref{A:capacity}: regressing 
capacity on the network variables using garage-level averages across 61 
garages with available data yields an adjusted $R^2$ of 0.008, and neither 
\textit{MindurOwnGar} nor \textit{CompGar10Min} are significantly 
correlated with capacity.}

A remaining concern is whether unobserved garage-level characteristics could 
bias the network variable estimates. The identifying assumption is that, 
conditional on the observed network variables, any remaining unobserved 
heterogeneity across garage locations is uncorrelated with the 
operator-specific regressors.\footnote{This assumption is non-trivial: 
the historical bus depots privatized in 1994 were not randomly located, 
potentially creating a correlation between unobserved site quality and 
the observed network variables. However, the severe building constraints 
in central London mean that garages are highly immobile, so that unobserved 
site characteristics are largely absorbed by the stable geographic features 
of the network directly captured by the included variables. The identifying 
assumption is therefore that conditional on route density, competitor 
proximity, and own-network centrality, no systematic unobserved 
heterogeneity remains.}

We verify empirically that the garage value function $\tilde{\Pi}(\cdot)$ satisfies supermodularity, confirming that the maintained assumptions for existence of a pure-strategy Nash equilibrium are consistent with the data ($\hat{\theta}_5 \leq 0$ and $\hat{\theta}_6 \leq 0$ as documented in Table~\ref{T:strucestim}). Our approach exploits the exclusivity of garage use and the temporal structure of our data by estimating on observed ownership \emph{transitions} rather than the full cross-sectional network, varying the number of garages assumed to choose operators simultaneously as a robustness check on equilibrium selection.


\subsubsection{Sequential discrete choice}\label{sec:seq-choice}
We now re-introduce the time dimension and use $t$ subscripts where variables 
are time-varying. In the vast majority of cases only one garage changes 
ownership at a given changedate (Table~\ref{T:garagechangedates}); we 
therefore first estimate the model under a sequential discrete choice 
framework with myopic decision-makers, treating each ownership change as an 
independent event. Cases where multiple garages change owner on the same 
changedate are treated as sequential draws from the same choice set, ordered 
as observed in the data. The observed garage ownership vector at each point 
in time is interpreted as a Nash equilibrium of the garage-chooses-operator 
game: conditional on the rest of the network remaining fixed, no garage has 
a profitable unilateral deviation to a different operator. This is the natural 
equilibrium concept for a static game, and is satisfied by construction when 
each garage is assigned to its highest-value operator given the observed 
assignments of all other garages.\footnote{In other words, when garage $j$ selects an operator at time 
$t$ according to \eqref{eq:garagechoice}, it conditions on the observed state 
of the network at that point ($\Gamma_t \equiv \mathbf{a}_{-jt}$).}

Adding the distributional assumption $\epsilon_{ijt} \overset{i.i.d.}{\sim} 
F_{\epsilon}(0,\sigma)$,\footnote{The location normalization $\mu = 0$ is without 
loss of generality: since $\epsilon_{ijt}$ is drawn i.i.d.\ from 
$F_{\epsilon}(\mu,\sigma)$ for all $i \in \mathcal{N}_t$ including the vacant 
option $i=0$, a common location $\mu$ cancels from every utility difference that 
determines choices and is therefore unidentified. This is distinct from the normalization 
of the vacant option's deterministic utility to zero, which pins down the level of 
$f(\cdot)$ and identifies the non-vacant intercept in $X$.} the choice 
probability for garage $j$ selecting operator $i$ at time $t$ is
%
\begin{align}\label{eq:choiceprob-general}
P_{ijt}(\theta) 
&= \Pr\!\Big(f(X_{ijt},G_{ijt};\theta) + \epsilon_{ijt} 
  \geq f(X_{hjt},G_{hjt};\theta) + \epsilon_{hjt} 
  \quad \forall\, h \in \mathcal{N}_t\Big) \\
&= \int \mathbf{1}\!\Big[f(X_{ijt},G_{ijt};\theta) + \epsilon_{ijt} 
  \geq f(X_{hjt},G_{hjt};\theta) + \epsilon_{hjt} 
  \quad \forall\, h \in \mathcal{N}_t\Big]\, dF_{\epsilon} \notag
\end{align}
where $\mathcal{N}_t$ is the set of active operators with at least one garage in the 
market at time $t$, always including the vacant choice $i=0$. By standard arguments, for any fixed $(X_t, G_t, \theta)$, the utility-maximising 
choice is unique with probability one, %\footnote{A tie between alternatives $i$ and $h$ requires 
%the error difference $\epsilon_{ijt} - \epsilon_{hjt}$ to equal the fixed 
%scalar $d_{ijht} \equiv f(X_{hjt},G_{hjt};\theta) - f(X_{ijt},G_{ijt};\theta)$. 
%Under any atomless $F_{\epsilon}$, the difference $\epsilon_{ijt} - 
%\epsilon_{hjt}$ has a continuous distribution, so 
%$\Pr(\epsilon_{ijt} - \epsilon_{hjt} = d_{ijht}) = 0$ for each $h \neq i$. 
%Since $|\mathcal{N}_t| < \infty$, a finite union of probability-zero events has probability 
%zero, and the maximiser is therefore unique with probability one.}
establishing uniqueness of garage-operator matches nonparametrically under mild conditions on $F_{\epsilon}$. When further specifying $F_{\epsilon}$ to be the standardized Type~I extreme value distribution --- which satisfies the conditions for uniqueness by being continuous 
and hence atomless ---, as we do in estimation, the integral in \eqref{eq:choiceprob-general} reduces to the 
\citet{mcfadden1984econometric} closed form:
\begin{align}\label{eq:entryprob}
P_{ijt}(\theta) = \frac{\exp\!\Big(f(X_{ijt},G_{ijt};\theta)\Big)}
{\sum_{h \in \mathcal{N}_t} \exp\!\Big(f(X_{hjt},G_{hjt};\theta)\Big)}
\end{align}
The logit functional form 
sets the scale parameter of $\epsilon$ to one, fixing the implied error variance 
at $\pi^2/6$, so estimated coefficients are scaled by $\sqrt{\pi^2/6}/\sigma$ relative 
to the true utility coefficients, where $\sigma$ is the standard deviation of the true 
unexplained component. The finite-dimensional parameter vector $\theta$ is estimated by 
maximum likelihood, using the observed identity of the operator entering each garage at 
each changedate and the network variables at that date.

The estimation strategy requires that the event of an ownership change is 
independent of $f(\cdot)$ conditional on the observables. This would be violated if firms make forward-looking entry 
decisions, for instance entering a location even at a short-run loss to 
preempt future competitor expansion in, e.g., \citep{goolsbee2008incumbents, zheng2016spatial}. While we account for a firm's tenure 
at a garage through the incumbency variable, abstracting from such dynamics 
is a limitation of the static framework, common to the entry literature more 
broadly \citep{Jia2008, ellickson2013estimating, nishida2015estimating}.\footnote{
\cite{Holmes2011} estimates a dynamic game with agglomeration effects for 
a single chain-store rollout.}

\subsubsection{Sequential bundled choice}\label{sec:seq-bundle}
As a robustness check that relaxes the purely static timing assumption, we 
also estimate the model based on the utility of \emph{bundles} of two 
consecutive ownership changes. The bundle framework is the natural first 
extension toward dynamics: by treating two sequential garage transactions as 
a joint decision, the entry of the first operator internalizes its 
consequences for the subsequent transaction, mimicking a one-step-ahead 
lookahead. Extending this to three or more steps is computationally 
prohibitive---the number of bundle combinations grows from 1,841 under 
sequential choice to 33,513 under two-step bundling---so we treat the 
two-step bundle as a tractable approximation. Comparing $\theta$ and 
$\theta^B$ serves as a diagnostic: if the spatial rent estimates are stable 
across the two specifications, this supports the static framework as a valid 
simplification; it would indicate longer planning horizons.

Let $B=\{(i_1,j_1),(i_2,j_2)\}$ denote the bundle in which operator $i_1$ 
enters garage $j_1$ and, in the subsequent transaction, operator $i_2$ enters 
garage $j_2$. Here $i_1$ and $i_2$ may be the same operator entering different 
garages in sequence; $j_1$ and $j_2$ may coincide only if $i_1 \neq i_2$, 
since exclusivity precludes the same operator entering the same garage twice 
consecutively. %When three or more garages change hands on the same changedate, 
%we form bundles from the first two and treat any additional changes as 
%independent sequential events.

The value of bundle $B$ is the average of the deterministic utility when $i_1$ 
enters $j_1$ and the deterministic utility when $i_2$ enters $j_2$ 
\emph{conditional on $i_1$ having entered $j_1$}, plus a bundle-specific 
unobservable $\epsilon^B \sim^{i.i.d.} F_{\epsilon^B}(0,\sigma^B)$:
\begin{align}\label{eq:garagevalue_detail_bundle}
\tilde{\Pi}^B(\Gamma) =\; &\theta_1^B \big( \text{YearsInc}_{j_1}+ 
  \text{YearsInc}_{j_2}\big)/2 \notag \\
+\; &\theta_2^B \big(\text{Nonvacant}_{i_1} + \text{Nonvacant}_{i_2}\big)/2 \notag \\
+\; &\theta_3^B \big(\text{log(NrGarages)}_{i_1} + \text{log(NrGarages)}_{i_2}(i_1)\big)/2 \notag \\
+\; &\theta_4^B \big(\text{RouteStarts10Min}_{j_1} + 
    \text{RouteStarts10Min}_{j_2}\big)/2 \notag \\
+\; &\theta_5^B \big(\text{MindurOwnGar}_{i_1 j_1} + 
    \text{MindurOwnGar}_{i_2 j_2}(i_1 j_1)\big)/2 \notag \\
+\; &\theta_6^B \big(\text{CompGar10Min}_{i_1 j_1} + 
    \text{CompGar10Min}_{i_2 j_2}(i_1 j_1)\big)/2 + \epsilon^B
\end{align}
Three network variables are evaluated conditional on $i_1$ entering $j_1$: 
$\textit{MindurOwnGar}_{i_2 j_2}(i_1 j_1)$ is the minimum drive time between 
$j_2$ and any garage of operator $i_2$ given $i_1$'s entry; 
$\textit{CompGar10Min}_{i_2 j_2}(i_1 j_1)$ is the number of competing garages 
within 10 minutes of $j_2$ for operator $i_2$ given $i_1$'s entry; and 
$\textit{log(NrGarages)}_{i_2}(i_1)$ is the natural logarithm of the number of garages of operator $i_2$ 
other than $j_2$, given $i_1$'s entry. The $\textit{Nonvacant}$ terms enter 
additively and are set to zero for the outside option $i=0$, so the bundle 
intercept equals $\theta_2^B$ when both garages are assigned to active 
operators, $\theta_2^B/ 2$ when one is assigned to an active operator and the 
other to the market, and zero when both are vacant---the outside option bundle. 
As above, all remaining variables are set to zero for the outside option bundle. The bundle choice 
probability is
\begin{align}\label{eq:entryprob-bundle}
P_{Bt}(\theta^B) = \frac{\exp\!\Big(f(X_{B},\Gamma_{B};\theta^B)\Big)}
{\sum_{c \in \mathcal{B}_t} \exp\!\Big(f(X_{c},\Gamma_{c};\theta^B)\Big)}
\end{align}
where $\mathcal{B}_t$ collects all possible bundles $(i_1,i_2) \in 
\mathcal{N}_t \times \mathcal{N}_t$ (with replacement), with $j_1$ fixed 
as the garage changing ownership at $t$ and $j_2$ as the next garage to 
transact. 

\florian{here we don't agree. a logit model with 2000 choice situations and one with 30000 has \emph{the same error variance}. What changes is the logsum (the expected value of the represented market, much higher in the 30K case). Again, comparing ratios of coefs may still make sense because our explanatory variables are different, but the argument cannot come from the error variance.}
Even if firms make only single-garage decisions throughout the sample, the bundle 
estimates $\theta^B$ may differ from $\theta$ in levels due to differences 
in the implied error variance $\sigma \neq \sigma^B$, which scales the population coefficients. % by 
%$\sqrt{\pi^2/6}/\sigma$ with $\sigma$ specific to each model. 
Since this 
scaling is common to all coefficients within a model, ratios of coefficients 
--- and hence the implied trade-offs between spatial rent determinants --- 
are scale-invariant and directly comparable across the two specifications. 
We therefore assess robustness by comparing these implied trade-offs rather 
than coefficient levels, using the same marginal rate of substitution 
between own and competitor garage proximity that is our main object of 
interest in Section~\ref{S:Results}.


\begin{table}[!tb]
% \centering
\caption{Summary statistics, location choice model}
\label{T:logitvars}
\begin{adjustbox}{max width = \textwidth, center}
	\input{\tables/logit-data-summary.tex}
\end{adjustbox}
\newline
\fnote{Summary statistics for variables used in the garage-operator utility model in equation \eqref{eq:garagevalue_detail} (top panel), for additional variables including alternative distance measures (middle panel), and for other characterizations of the discrete choice model (bottom panel). The top two panels are based on all 1,649 observations (ijt) in the estimation sample where the garage is non-vacant (e.g., applying to all active operators, $ i\neq 0$) and the summary statistics for the variables ``Alternative is selected" and ``Garage is not vacant" in the bottom panel are based on the full estimation sample of 1,841 observations, including $i=0$. The variables capturing economies of density are 0 when the operator has only one garage.}
\end{table}


\section{Estimation results}\label{S:Results}
This section presents the estimation results of the \emph{garages choose operators} 
model introduced in Section~\ref{sec:seq-choice}, under the assumption that 
when a garage changes ownership, the entrant is the operator with the highest 
utility for that garage among all active operators at that time.

\subsection{Baseline estimates}
The main estimation results are presented in Table~\ref{T:strucestim}. We 
discuss each component of the garage utility function in turn based on the 
baseline sequential estimates, before turning to a comparison with the bundle 
estimates in Section~\ref{sec:bundle-results}.

\paragraph{Incumbency benefit.} The coefficient on $\textit{YearsInc}$---the 
number of years the incumbent has been in the garage---is positive and 
significant ($\hat{\theta}_1 = 0.121$, s.e. $0.040$, $p<0.01$). The linear 
specification implies a constant rate of accumulation of incumbency advantage 
over time: each additional year of tenure increases the incumbent's utility 
for the garage at a constant rate, reflecting the gradual build-up of 
garage-specific investments and operational logistics that are costly for a 
new operator to replicate.\footnote{An alternative parameterization using 
$\exp(-\textit{YearsInc}_j)$ is arguably more naturally motivated by the five-year 
contract structure: if an operator enters a garage to serve one or more 
contracts immediately after winning, the incumbency advantage would be 
highest at entry and decay toward zero over the contract duration. However, 
the robustness checks in Section~\ref{SS:robust-inc} show that the linear 
specification fits the data better. This may reflect the staggered nature 
of route contracts---whereby an incumbent simultaneously holds contracts at 
different stages of their duration---so that the garage retains residual 
value throughout the tenure rather than only at the start.}

\paragraph{Alternative market usage.} Active operators have significantly 
higher mean utility for a garage than the outside option ($\hat{\theta}_2 
= 1.886$, s.e. $0.529$, $p<0.001$), confirming that bus operators value 
garages substantially more than their alternative use value, consistent 
with the garage being a productive asset specifically tied to TfL route 
procurement.


\paragraph{Scale effects.} The coefficient on $\log(\textit{NrGarages})$---the 
log of the total number of garages operated by the firm---is negative and 
highly significant ($\hat{\theta}_3 = -1.135$, s.e. $0.124$, $p<0.001$). 
The logarithmic specification implies diminishing returns to garage expansion: 
each additional garage contributes less to operator utility as the network 
grows larger. Conditional on network centrality and competitor proximity, 
larger operators derive less marginal utility from an additional garage, 
consistent with firms having fewer uncovered routes to serve from a new 
garage as their network becomes denser.\footnote{The logarithmic specification fits the data better than a linear 
specification in levels, as shown in the robustness checks of 
Section~\ref{SS:robust-inc}.}


\paragraph{Dead miles.} The coefficient on $\textit{RouteStarts10Min}$---the 
number of route endpoints within a 10-minute drive---is positive as expected 
but statistically insignificant in the baseline specification 
($\hat{\theta}_4 = 0.030$, s.e. $0.028$). This is consistent with the 
reduced-form evidence in Section~\ref{S:TfL}, where the dead miles 
coefficient was similarly attenuated once strategic network controls were 
included. In the baseline, the dead miles variable may be partially collinear 
with the non-vacancy intercept, both of which partly reflect garage centrality 
in the route network. The bundle model, by internalizing the network 
consequences of sequential entry, provides better identification of this 
component: the bundle estimate is large and highly significant 
($\hat{\theta}_4^B = 0.161$, s.e. $0.023$, $p<0.001$). We therefore treat the 
baseline estimate as a conservative lower bound on the true dead miles 
effect while noting that both specifications agree on the sign and that the 
dead miles coefficient does not affect the robustness of the other spatial 
rent estimates, which are stable across both specifications.

\paragraph{Economies of density.} The coefficient on 
$\textit{MindurOwnGar}$---the minimum drive time to the operator's nearest 
own garage---is negative and highly significant ($\hat{\theta}_5 = -0.091$, 
s.e. $0.010$, $p<0.001$). Garages are more valuable when they are closer to 
the operator's existing network, confirming the presence of agglomeration 
benefits. These provide a cost-based rationale for the observed spatial 
clustering of operators in certain parts of London.

\paragraph{Local monopoly rents.} The coefficient on 
$\textit{CompGar10Min}$---the number of competitor garages within a 
10-minute drive---is negative and significant ($\hat{\theta}_6 = -0.334$, 
s.e. $0.110$, $p<0.05$). Garages are less valuable when surrounded by 
competitor garages, consistent with the local monopoly rent mechanism: 
proximity of rivals compresses equilibrium mark-ups in route auctions, 
reducing the expected profit from operating routes nearby. To quantify the relative importance of economies of density and local 
monopoly rents, we compare the ratio $|\hat{\theta}_6|/|\hat{\theta}_5| 
= 3.7$, which gives the reduction in minimum own-garage drive time (in 
minutes) that offsets one additional competitor garage within 10 minutes. 



\begin{table}[!tb]
\centering
\caption{Structural estimates of garage-operator utility parameters}\label{T:strucestim}
\begin{adjustbox}{max width = 0.9\textwidth}
	\input{\tables/L-paper-preferred-model.tex}
\end{adjustbox}\fnote{Statistical significance: + p $<$ 0.1, * p $<$ 0.05, ** p $<$ 0.01, *** p $<$ 0.001. The column ``baseline" contains the structural estimates from the sequential single-move discrete choice model based on equation \eqref{eq:garagevalue_detail}. The column ``bundle"  contains the model estimates when bundles of two garages and operators are chosen sequentially, based on the utility function in equation \eqref{eq:garagevalue_detail_bundle}.} \end{table}

\begin{table}[!tb]
\centering
\caption{Testing structural estimate equality across baseline and bundle specifications}\label{T:deltamethod}
\begin{adjustbox}{max width = \textwidth}
	\input{\tables/delta-method.tex}\end{adjustbox}
\fnote{Ratios between the coefficients on the listed variable 
to the coefficient on $\textit{CompGar10Min}$ in the relevant specifications reported in Table~\ref{T:strucestim}. 
Ratios are scale-invariant and directly comparable across specifications 
despite differences in implicit error variance $\sigma$. The $p$-value is 
from a two-sided $z$-test of equality of the two ratios, as defined in 
\eqref{eq:deltamethod}, with standard errors obtained via the delta method \citep{greene2012econometric}
applied to the respective maximum likelihood covariance matrices.}\end{table}


\subsection{Bundle choice estimates}\label{sec:bundle-results}
As discussed in Section~\ref{sec:seq-bundle}, the bundle model allows 
operators to internalize the network consequences of sequential entry, 
providing a robustness check on the static baseline. Since the two 
specifications differ in their implicit error variance $\sigma$, coefficient 
levels are not directly comparable across columns. We therefore assess 
robustness by comparing coefficient ratios relative to $\hat{\theta}_6$ 
(the competitor garage coefficient), as such ratios are scale-invariant and directly 
interpretable as trade-offs between garage value determinants.

To formally test whether these ratios differ across specifications, we apply 
the delta method to obtain standard errors  \citep{greene2012econometric} for each ratio within each model 
and construct a $z$-test of equality exploiting the independence of the two 
estimators:
\begin{equation}\label{eq:deltamethod}
z_k = \frac{\hat{\theta}_k/\hat{\theta}_6 - \hat{\theta}_k^B/\hat{\theta}_6^B}
{\sqrt{\widehat{\text{Var}}(\hat{\theta}_k/\hat{\theta}_6) + 
\widehat{\text{Var}}(\hat{\theta}_k^B/\hat{\theta}_6^B)}}
\end{equation}
for $k=1:5$. Variances are obtained via the delta method applied to the respective 
maximum likelihood covariance matrices. The results are reported in 
Table~\ref{T:deltamethod}.

The test fails to reject equality for four of the five coefficients, 
including the key trade-off between economies of density and local monopoly 
rents ($\hat{\theta}_5/\hat{\theta}_6$: baseline $0.273$, bundle $0.282$, 
$p=0.945$). This is the primary object of interest---the number of minutes 
of own-garage proximity that offsets one additional competitor garage---and 
its near-identical value across specifications provides strong support for 
the baseline estimates. Scale effects and the incumbency benefit are likewise 
statistically indistinguishable across models ($p=0.846$ and $p=0.688$ 
respectively). The one ratio that does differ significantly across 
specifications is that of dead miles ($p=0.019$): the bundle estimate of 
$\textit{RouteStarts10Min}$ is large and significant (Table~\ref{T:strucestim}) 
while the baseline is not, so the ratio to \textit{CompGar10Min} is roughly 
four times larger in the bundle model. As discussed above, this reflects 
better identification of the dead miles component in the bundle model rather 
than instability in the spatial rent estimates. Importantly, this difference 
does not affect the robustness of the primary findings: local monopoly rents 
and economies of density---the main drivers of the counterfactual analysis---are stable across both specifications.

Taken together, these results support the static sequential framework as a valid 
simplification for the purpose of identifying the spatial rents in this environment. Extending the planning horizon from one to two steps does 
not materially alter the picture of how spatial positioning drives garage 
value in this market. At the same time, we cannot rule out that alternative 
bundling structures---for instance, non-sequential bundles or three-step 
lookaheads--would yield different estimates, particularly for the dead miles 
component where the two specifications do differ. What the comparison does 
establish is that a parsimonious static model without prices or explicit 
dynamics is able to recover stable and economically interpretable estimates 
of the two central forces---agglomeration benefits and local monopoly 
rents---that theory predicts should govern location choice in markets with 
spatial competition. 

\subsection{Model fit} 
Multinomial choice models do not admit a single summary fit statistic 
analogous to $R^2$. We therefore assess fit using the mean Area Under 
the ROC Curve (AUC), averaged across operators using a one-vs-all 
classification scheme and weighted by operator frequency in the 
sample.\footnote{The AUC measures the probability that the model assigns 
a higher predicted probability to the true entrant than to a randomly 
drawn non-entrant, ranging from 0.5 (random assignment) to 1 (perfect 
discrimination). In the multinomial setting we compute one ROC curve per 
operator under a one-vs-all scheme and report the frequency-weighted 
average following \cite{provost2003tree}.} The baseline model achieves 
a mean AUC of 0.868 and a McFadden pseudo-$R^2$ of 0.554; the bundle 
model achieves an AUC of 0.875 with a pseudo-$R^2$ of 0.422, the lower 
value reflecting the much larger choice set rather than worse 
fit.\footnote{McFadden pseudo-$R^2$ values of 0.2--0.4 are conventionally 
considered to indicate excellent fit, substantially exceeding the threshold 
for OLS regressions of the same magnitude \citep{domencich1975urban}. The 
baseline value of 0.554 is therefore exceptional for a model with only five 
network variables.} The AUC metric, which is robust to the number of 
alternatives, confirms that both specifications have strong discriminatory 
power and that the key spatial rent parameters are stable across estimation 
approaches. The model correctly identifies the highest-utility operator in 
50.5 percent of all choice situations, out of an average of 10.75 alternatives 
per situation---a substantial improvement over random assignment.

Figure~\ref{F:frequencies} compares in-sample predicted and observed entry 
frequencies, computed as the model-implied average probability that each 
operator is selected across all choice situations against the observed share 
of ownership changes in which that operator entered a garage. The model 
tracks individual operator entry patterns closely across the distribution. 
The main discrepancy is an over-prediction of vacancy: the model assigns 
garages to the outside option in 38.5 percent of cases against 26.0 percent 
in the data. This reflects the difficulty of capturing the timing of garage closures and 
developments, which often occur for reasons exogenous to bus operations, and 
supports the modelling choice to treat the set of garages as the fixed outer 
set of ever-occupied plots rather than modelling entry and exit of garage 
infrastructure explicitly.

A second discrepancy highlighted in Figure~\ref{F:frequencies} is that the model under-predicts entry probabilities 
for Arriva and Go-Ahead. This follows mechanically from the estimated 
diseconomies of scale in $\log(\textit{NrGarages})$ combined with these 
operators' large garage holdings (Figure~\ref{fig:garages-TS-byoperator}): 
conditional on network centrality and competitor proximity, the model 
assigns low marginal value to further expansion by already large operators. 
Appendix~\ref{A:collusion} explores a complementary explanation, motivated 
by the UK Competition Commission's concerns about geographic market 
segregation in the London bus market: whether unrealized garage transactions 
that the model predicts to be profitable are consistent with a collusive 
market-sharing interpretation. The most robust finding is that implied 
punishment values (under this interpretation) are systematically larger for garages that are more 
spatially isolated from competitor networks, directly connecting the pattern 
of unrealized transactions to the local monopoly rent estimates in 
Table~\ref{T:strucestim}. Results relating punishment values to the number 
of active firms are sensitive to specification; we therefore stop short of 
claiming unequivocal evidence of collusion and refer the reader to 
Appendix~\ref{A:collusion} for full details.\footnote{It is worth recalling that the spatial rent estimates are 
identified entirely from observed ownership transitions, without using 
procurement bid data, and are agnostic about competitive conduct at the 
auction stage. The only assumption required is that transactions that do 
take place allocate garages to the operator that values them most---the 
model makes no use of information about transactions that did not occur.} Here, we simply conclude that the spatial rent model fits the observed network changes remarkably well given its parsimony, without explaining specific operator-specific deviations from 
its predictions.

\subsection{Alternative specifications}\label{S:Robustness}
We assess the robustness of the baseline estimates to alternative 
specifications of the distance measures, the incumbency variable, and 
the scale effect.

\subsubsection{Alternative distance measures}\label{SS:robust-dist}

Appendix Table~\ref{T:robustness_dist} reports estimates under alternative 
specifications of the three distance components; we discuss each substitution 
in turn. Replacing \textit{RouteStarts10Min} with stopover-based dead mile 
measures leaves all spatial rent estimates unchanged and the stopover 
coefficients themselves are insignificant, confirming that the dead miles 
attenuation in the baseline is not specific to the start-stop measure. 
Replacing \textit{MindurOwnGar} with the average drive time to own garages 
yields a stable local monopoly rent coefficient but a much larger trade-off 
ratio (9.9 vs. 3.7 in the baseline) and lower fit (AUC 0.844 vs. 0.868), 
reflecting the different scaling of the average relative to the minimum and 
confirming the minimum as the preferred summary of agglomeration benefits. 
Replacing \textit{CompGar10Min} with continuous drive-time measures yields 
insignificant competitor distance coefficients, suggesting that the count 
of nearby rivals is more informative than their average distance. The 
baseline specification compares favorably on all fit metrics and the core 
spatial rent estimates are robust to the choice of distance measure.

\subsubsection{Alternative incumbency and scale 
specifications}\label{SS:robust-inc}

Appendix Table~\ref{T:robustness_other} reports estimates under alternative 
specifications of the incumbency and scale variables. Replacing the linear 
incumbency term with the exponential form $\exp(-\textit{YearsInc})$ yields 
an insignificant coefficient and lower fit, confirming the preference for 
the linear specification. Omitting incumbency entirely leaves the spatial 
rent estimates and fit virtually unchanged, confirming that the incumbency 
specification does not drive the primary findings. Replacing 
$\log(\textit{NrGarages})$ with \textit{NrGarages} in levels reduces fit 
noticeably while leaving the trade-off ratio stable; omitting the scale 
variable entirely reduces fit further without materially affecting the 
spatial rent estimates. Omitting the \textit{Nonvacant} intercept restores 
significance of the dead miles coefficient, further supporting the 
collinearity argument in Section~\ref{S:Results}. Throughout, the 
trade-off ratio between economies of density and local monopoly rents 
ranges from 3.1 to 4.7, and the core finding that these two forces are 
the primary determinants of garage value is robust across all alternative 
specifications considered.


\section{Counterfactuals: implications for urban planning}\label{S:Counterfactuals}
The structural estimates link garage value directly to spatial positioning 
in the network of routes and competing garages. We exploit this to assess 
two urban planning questions: where in London would new garage infrastructure 
generate the most private value to operators, and where are capacity constraints on existing 
garages most binding? The first question is addressed descriptively using 
the estimated spatial rent surface across candidate locations; the second 
requires the bus-garage assignment equation, which we validate against 
observed tender data before applying it to the capacity expansion 
counterfactuals. 

\subsection{Network quality of candidate garage locations}\label{SS:UP3}
For each point $\ell$ on a 500-metre resolution grid covering Greater London, 
we compute the three spatial rent components of the reduced-form garage utility 
function \eqref{eq:garagevalue_detail} as if a new garage were located at 
$\ell$.\footnote{These are the number of route endpoints within a 10-minute drive, the minimum drive time from $\ell$ to the 
nearest existing garage of each operator $i$, 
and the number of competitor garages within a 10-minute drive of $\ell$. Drive times are computed using the OSRM 
routing engine on the London road network.} For each candidate location $\ell$ 
and each operator $i$, we then evaluate $\hat{\Pi}_{i\ell}(\Gamma_{-\ell})$ using the 
baseline structural estimates, excluding any garages situated inside the 500 m$^2$ area $\ell$ but holding the rest of the 2019 garage-operator 
network $\Gamma$ fixed. We also determine the highest-value operator at each location, as \footnote{In what follows, $\hat{\Pi}$ denotes the estimated deterministic garage utility given our baseline estimates.}
$$i^*(\ell) = \argmax_i \hat{\Pi}_{i\ell}(\Gamma_{-\ell})$$
This exercise is purely 
descriptive: it identifies where in London the marginal value of additional 
operating capacity is largest for each operator, based on the spatial 
positioning of candidate locations relative to routes and existing garages. 
It does not use the bus-garage assignment in equation \eqref{eq:assignment_buses} and abstracts from garage 
capacity---questions we take up in Sections~\ref{SS:urbanpolicy1}--\ref{SS:likecompanion}.\footnote{Incorporating bus-garage assignment can address issues like the possibility that a garage in an otherwise 
unserved area could generate stand-alone value by accessing routes not 
currently viable from existing garages.} Figure~\ref{F:urban3_operator} presents the two complementary views of the 
spatial rent surface. Panel (a) shows the overall \emph{private} value of 
releasing land for garage development at each candidate location, 
independently of which operator would benefit. Rather than a clear geographic 
gradient, the surface reveals many dispersed high-value pockets scattered 
across London, reflecting that the 2019 garage network is broadly 
well-positioned relative to the route network---each operator has already 
claimed the most valuable nearby locations, leaving isolated pockets of 
residual private value where the network remains sparse. The northeast is 
the main exception, where a larger contiguous area of higher private value 
suggests that garage coverage is genuinely underprovided relative to the 
density of the route network there. 

Panel (b) maps $i^*(\ell)$, the highest-value operator at each location. 
The spatial segmentation is clean and contiguous: Stagecoach dominates large areas in the south and south-east, Abellio the north-west, and smaller operators 
occupy distinct central territories. Most existing garages lie within their 
own operator's preferred region, confirming that the current network broadly 
reflects spatial rent optimization. %The main exception is the large northern 
%region attributed to HCT, one of the smallest operators, rather than Arriva, 
%whose garages are closest to this area. This directly reflects the 
%diseconomies of scale in the number of garages: large operators derive less 
%marginal utility from additional capacity, so the map should be read as 
%showing where extra capacity would benefit each operator most given their 
%current network, not where absolute garage value is highest.
At a high level, the results reveals where operators would 
have the strongest incentive to lobby for the release of industrial land or 
planning permissions for new garage development, with the northeast emerging 
as the area where private returns to new garage infrastructure would be 
largest.\footnote{The underlying geography is relevant: this area is situated in the middle of large stretches of \href{https://apps.london.gov.uk/planning/}{designated open space}, protected against development.}
%the garage placed at the center of the high value area is (\florian{find garage name}), in the Loughton, bordering to the left with the large Epping Forest - a natural reserve protected from development.

\begin{figure}[!tb]
\caption{Network quality surface for new garage locations\label{F:urban3_operator}}
\begin{subfigure}[b]{0.49\textwidth}
    \includegraphics[width=\textwidth]{\plots/urban3-value-heatmap.pdf}
    \caption{Location value surface}
\end{subfigure}
\hfill
\begin{subfigure}[b]{0.49\textwidth}
    \includegraphics[width=\textwidth]{\plots/urban3-operator-heatmap.pdf}
    \caption{Optimal entering operator}
\end{subfigure}
\fnote{Panel (a) shows the highest 
predicted utility across all operators on 
a 500-metre grid, reflecting the overall value of releasing land for garage 
development at that location. Panel (b) 
shows the operator that 
would benefit most from a new garage at each location. Dots mark 
existing 2019 garages. Both panels are evaluated using the 2019 
garage-operator network $\Gamma$ and reflect pure network position value, 
abstracting from garage capacity and bus-garage assignment.}
\end{figure}

\subsection{Private returns to garage capacity expansion}\label{SS:UP1}
\label{SS:urbanpolicy1}
This section quantifies the returns to expanding capacity at existing garages with 50 buses, which is approximately half of the mean (and median) capacity in the data and asks how much private profit the
owning operator would gain. The answer combines two components, each capturing
a distinct channel through which additional bus capacity creates value. 

The first channel is the \emph{structural term}. The structural model assigns
each garage a deterministic valuation $\hat{\Pi}_{ij}$ that reflects the
garage's position in the operator's network: its proximity to own and
competitor garages, the density of nearby route starts, and the operator's
incumbency. Under the assumption of constant returns to scale in capacity
(motivated in Appendix~\ref{A:capacity}), dividing $\hat{\Pi}_{ij}$ by the
number of buses currently stationed at the garage, $B_{jt}$, yields a
per-bus valuation $\hat{\pi}_{ij}^{bus} = \hat{\Pi}_{ij}/B_{jt}$. Multiplying
by the number of additional buses drawn into the garage and converting to
monetary units via the implied TfL auction revenue per utility unit
$\kappa = \text{\textsterling}1.23\text{m}$, gives the structural component of
the return to expansion.

The second channel is the \emph{dead miles correction}. The structural model
is estimated at the garage level and does not observe which specific routes are
assigned to which garage; the variable $\mathit{RouteStarts10Min}$ captures
route proximity only through a binary ten-minute threshold. When capacity
expands, the assignment technology of equation~\eqref{eq:assignment_buses}
predicts which additional routes get stationed at the garage: buses are added
in order of route proximity until 50 extra buses are absorbed. These marginal
routes need not have the same dead-mile profile as the routes already assigned
to the garage. If marginal routes are closer than the baseline average, the
operator saves on empty driving; if they are further, costs rise. The
dead miles correction accounts for this difference by comparing the
dead-mile cost of each marginal route $m'$ to the PVR-weighted average dead
miles of routes already stationed at the garage, $\overline{DM}_j$, and
scaling by the estimated dead-mile penalty
$\hat{\beta} = -\text{\textsterling}19{,}000$ per additional dead mile driven
(from Section~\ref{S:TfL}). Since $\hat\beta < 0$, the correction is negative
whenever marginal routes are further away than the baseline average
($(DM_j(m') - \overline{DM}_j) > 0$), and positive when they are closer.

Combined, the private return to operator $i$ from adding
50 buses to garage $j$ is:
%
\begin{equation}\label{eq:UP1}
  \Delta\hat{\Pi}_{ij}(+50\text{ buses}) =
  \underbrace{%
    \sum_{m' \in \mathcal{M}_j}
    \Delta b_j(m') \cdot \hat{\pi}_{ij}^{bus} \cdot \kappa
  }_{\text{structural term}}
  +
  \underbrace{%
    \hat{\beta} \cdot \sum_{m' \in \mathcal{M}_j} \Delta b_j(m')
    \cdot \bigl(DM_j(m') - \overline{DM}_j\bigr)
  }_{\text{dead miles correction}}
\end{equation}
%
where $\mathcal{M}_j$ is the set of marginal routes drawn in beyond the
current assignment, ordered by proximity according
to~\eqref{eq:assignment_buses}; $\Delta b_j(m')$ is the number of additional
buses from route $m'$ stationed in the garage (equal to the full PVR of $m'$
for all but the last marginal route, which is truncated so that the total
increment is exactly 50 buses); $\overline{DM}_j$ is the PVR-weighted average
dead-mile cost of routes already stationed at the garage at baseline; and all
other terms are as defined above. Further details and justification of these
measures are provided in Appendix~\ref{A:Capacity_expansion_details}.


%The empirical validity of the assignment technology 
%underlying $\mathcal{M}_j$ is supported by the validation in 
%Appendix~\ref{A:assignment_validation}, which recovers the correct operating 
%garage for 46 percent of routes exactly and 62--70 percent within a 
%two-minute dead-mileage tolerance.

 
Figure~\ref{F:term_decomposition} presents the joint distribution of the two
components of~\eqref{eq:UP1} across garages.\footnote{Full rankings are
reported in Appendix Tables~\ref{tab:cap-expansion-binding}
and~\ref{tab:cap-expansion-slack}.} The horizontal axis measures the
structural per-bus garage value in monetary units via $\kappa$, recentered
relative to the lowest-valued garage in the sample.\footnote{The garage with
the lowest $\hat{\pi}^{bus}$ is Grays, the easternmost garage in the sample,
located just beyond the Greater London boundary and owned by Arriva.} The
vertical axis measures the dead miles correction, negative when marginal
routes are further than the baseline average dead miles of currently assigned
routes and positive when they are closer.\footnote{The correlation between
the two terms is positive ($r = 0.161$): structural estimates capture that
garages closer to more routes have higher value regardless of where those
routes are currently assigned, while variation in the dead miles penalty
conditional on structural value reflects route-garage-specific empty driving
times that go beyond the binary 10-minute threshold in
\textit{RouteStarts10Min}.}
 
The positive dead miles correction for some garages identifies cases where
the observed allocation includes relatively distant routes that inflate
$\overline{DM}_j$, suggesting inefficiency in the baseline assignment
relative to proximity. For most garages, however, the correction is negative:
marginal routes are further away than the baseline average, underscoring the
increasing marginal cost of capacity expansion as closer routes are already
served. This implies that even absent binding land constraints, dead mile
costs alone are sufficient to make multiple smaller garages more profitable
than a single large depot for a profit-maximizing operator.
 
The figure highlights the two distinct routes to high expansion value. Erith
and Kangley Bridge Road rank near the top on the structural dimension despite
modest dead miles corrections; Waterloo reaches a similar total return
through a large positive dead miles correction, reflecting that its currently
assigned routes are more distant than those a capacity increment would draw
in. At the other end, Bow, Tottenham, and Northumberland Park are pulled to
the bottom of the distribution by strongly negative dead miles corrections:
the routes available at the margin are far from the depot, and this cost
dominates any structural gain.
 
To explain what drives the private returns to investment, we regress
$\Delta\hat{\Pi}_{ij}$ on the structural model covariates, capacity,
utilization, and the average baseline dead miles, with and without operator
fixed effects (Table~\ref{tab:cap-expansion-reg}). In line with the
discussion above, an important factor is the average dead miles of currently
assigned routes: garages serving more distant routes at baseline have lower
expansion value since their marginal routes will be even further
away.\footnote{The baseline dead miles measure is not included in the
structural model --- as it is only observed for garages in the final
allocation and not throughout the sample --- where route proximity is
captured by the binary \textit{RouteStarts10Min} threshold. The two empty
driving measures are collinear: conditional on baseline dead miles, garages
close to many route endpoints they do not currently serve face a larger dead
miles penalty on marginal expansion, which explains the negative coefficient
on \textit{RouteStarts10Min} in columns~(1)--(2). When baseline dead miles
are excluded (column~3), \textit{RouteStarts10Min} becomes insignificant
and $\log(\textit{NrGarages})$ and capacity become significant, reflecting
that operators with many large garages tend to be more distant from the route
network --- as illustrated by Arriva's Grays garage at the eastern
boundary.} Incumbency tenure is a robust positive driver across all
specifications, reflecting that firms stay longer in more profitable garages.
Operator fixed effects absorb substantial variation --- $R^2$ rises from
$0.876$ to $0.962$ --- reflecting that operator-specific network
configurations generate systematically different baseline expansion values,
ranging from \pounds$12.8$m for Tower Transit to \pounds$21.6$m for Arriva
after controlling for garage-level covariates, and consistent with firm
heterogeneity in purely network-based entry opportunities documented in
Section~\ref{SS:UP3}.
 
The key policy-relevant finding concerns competitor proximity. Conditional on
operator fixed effects, one additional competitor garage within a 10-minute
drive reduces the return to a 50-bus capacity expansion by \pounds$735{,}000$
on average --- approximately 3--6 percent of the operator-specific baseline
expansion value. This quantifies the wedge between private and social
returns: operators are incentivized to protect spatial isolation from
competitors, which generates local monopoly rents but imposes procurement
costs on the taxpayer, rather than to expand toward underserved parts of the
route network.
 
A central concern in our framework, noted in the introduction, is that
location-varying fixed costs are not identified. Garages in
competitor-sparse areas may simply be cheaper to expand because they tend to
lie outside high-land-value city-centre boroughs, conditional on closeness to
the route network. Columns~(4)--(6) of Table~\ref{tab:cap-expansion-reg}
address this by deflating $\Delta\hat{\Pi}$ by a borough-level house price
index, demeaned so that the average London borough takes the value 1 and
ranging from 0.83 to 1.15 across the garages in the sample. The adjustment
has a modest effect on the competitor-proximity estimate:
$-$\pounds$679{,}000$ in column~(5) versus $-$\pounds$735{,}000$ unadjusted
in column~(2).\footnote{After excluding garages in the periphery, operator
fixed effects reverse sign, and capacity expansion becomes more profitable
for larger operators. Besides sample selection, this is consistent with
different firms operating from systematically different parts of the city,
with operator identity partially proxying for location in high-HPI inner
boroughs in the nominal regressions.} Columns~(7)--(9) go further, deflating
by ward-level house prices to absorb finer-grained variation in land values
within boroughs. The competitor-proximity effect strengthens to
$-$\pounds$1{,}142{,}000$ in column~(8), suggesting that if anything the
borough-level adjustment understates the spatial-isolation incentive by
failing to capture within-borough land-value gradients. Arguably, residential
house prices are an imperfect proxy for industrial land values and
development costs. However, to the extent that they capture differences in
land values, city-centre density, and depot redevelopment costs, the
stability and strengthening of the competitor-proximity coefficient across
columns~(2), (5), and (8) minimizes the concern that the spatial-isolation
incentive is an artefact of peripherally located garages facing lower
unobserved fixed expansion costs.
 
We conclude that the socially optimal margin of investment would minimize
dead miles and operational costs; the privately optimal margin additionally
rewards competitor isolation. Regulation has value precisely because this
wedge is not internalized by profit-maximizing operators, as further explored
in Appendix~\ref{A:collusion}.

%$25.7 \times (\log 11 - \log 10) \approx \text{\textsterling}1\text{m}$, 
%Co$\text{\textsterling}34.6\text{m}$ 
%in column (2) implies that a hypothetical operator with a single garage and 
%average values of all other covariates would find capacity expansion highly 
%profitable---it is the accumulation of garages that erodes this value, 
%consistent with the structural model's prediction that market power is most 
%concentrated among operators with small, strategically positioned networks 
%rather than large ones. The $R^2$ of 0.36--0.45 confirms that the structural 
%model covariates account for a substantial share of the variation in 
%$\Delta\hat{\Pi}_{ij}$, supporting the use of \eqref{eq:UP1} as a coherent 
%summary of the returns to capacity expansion.

\begin{figure}[!tb]
\caption{Private returns to capacity expansion}\label{F:term_decomposition}
\centering
\includegraphics[width=0.8\textwidth]{\plots/urban1-term-correlation.pdf}
\fnote{Structural term ($\hat{\pi}_{ij}^{bus}$ horizontal axis) and dead miles correction 
(vertical axis) from returns to capacity expansion \eqref{eq:UP1}, for all 62 garages in the sample with observed capacity that are not vacant and have buses stationed in 2019. Garages are coloured by baseline capacity utilization and sized 
by capacity. Labeled garages rank among the three highest on the structural 
term, the two most extreme on the dead miles correction, or the four largest 
by capacity. Grays serves as a reference point with lowest $\hat{\pi}_{ij}^{bus}$.}
\end{figure}
%\caption{Urban policy 1 -- term1 on x-axis and term2 on y-axis\label{F:term_decomposition}}

\subsection{Welfare-maximizing global allocation}\label{SS:likecompanion}                                                                                              
%This section quantifies how much welfare could improve if the garage-operator                                                                                        
%assignment itself were optimized. Since TfL procures routes at cost plus                                                                                               
%mark-up, minimizing total expected operator profit $\sum_j \hat{\Pi}_{ij}(\Gamma)$
%is equivalent to minimizing total procurement payments by the taxpayer. We                                                                                             
%therefore find the assignment $\Gamma^*$ that achieves this. In \cite{marra-oswald} we show that unbundling garage                                                  
%ownership from route operation and central assignment---allowing a system operator to assign buses                                                                     
%to any garage regardless of ownership---has the potential to reduce dead miles by 14 percent. That paper abstracts from strategic garage interactions, which yields a linear assignment  
%problem solvable with a standard LP solver.                                                                                                                            
                                                                                                                                                                    
%Unlike that linear benchmark, the present problem, where garage utility depends on the full allocation through strategic network interactions, is a variant of the   
%quadratic assignment problem and NP-hard, in general: reassigning a single garage across two operators cascades through the objective of every other garage held by both affected operators 
%via three coupled, non-separable terms --- the garage count, the minimum distance to the operator's own network, 
%and the local competitor count. The feasible space of $8^{68} \approx 10^{61}$ candidate allocations of garages to operators in 2019 rules 
%out exhaustive search, motivating simulated annealing with multi-restart. \footnote{This is similar to the algorithm in \cite{kreindler2024optimal}, and exploits the
%fact that drive times between all garages and routes are precomputed and fixed, so only the network variables of garages within a 10-minute radius of each proposed
%reassignment need to be updated at each step. Simulated annealing is a stochastic local search algorithm that accepts welfare-improving swaps with probability one and
%welfare-reducing swaps with a positive probability that declines over time according to a cooling schedule, allowing escape from local optima early in the search. We
%run three restarts, each initialised from the best state found so far, and report the allocation achieving the highest welfare across all restarts.\florian{this needs to be fixed --- Florian I leave the fill in of the technical details to you, ok?}}

The capacity-expansion counterfactuals quantify the private--social wedge
arising from garage location choices. We now turn to simulations that estimate
the welfare gains from reassigning garages across operators, holding the
existing route and garage network fixed.
 
The structural estimates imply that the primary source of inefficiency is the
private incentive to maintain spatial isolation from competitors. Unlike the
cost-driven components of the profit function---zero-sum transfers between TfL
and the winning operator---local monopoly rents constitute a deadweight loss:
by locating away from competitors, operators sacrifice proximity to routes and
own garages, increasing dead miles and the associated externalities (pollution,
noise, congestion) borne by society.%
\footnote{A second source of inefficiency captured by these counterfactuals
stems from insufficient garage transactions, labelled ``ownership frictions''
in \citet{marra-oswald} and discussed further in Appendix~\ref{A:collusion}.
Land-use restrictions, analysed in Sections~\ref{SS:UP3} and~\ref{SS:UP1},
constitute a third source that we do not address here; garage locations and
sizes are held fixed throughout.}
A social planner therefore seeks to minimise local monopoly rents while holding
operational costs no higher than at the 2019 baseline.
 
Formally, the planner solves
%
\begin{align}
  \label{eq:planner}
  W_1 = \max_{\Gamma'} \; &
    -\sum_{i,k} \left[
      \hat{\Pi}_{ik}(\Gamma')
      + \hat{\theta}_1 \cdot \Delta\mathrm{YearsInc}_{k}(\Gamma')
    \right], \\
  \text{subject to:}\quad
    & \sum_{i,k} \mathrm{MinDurOwnGar}_{ik}(\Gamma')
      \leq \sum_{i,k} \mathrm{MinDurOwnGar}_{ik}(\Gamma),
      \label{eq:c1} \\
    & \sum_{k} \log(\mathrm{NrGarages})_{ik}(\Gamma')
      \leq \sum_{k} \log(\mathrm{NrGarages})_{ik}(\Gamma),
      \label{eq:c3} \\
    & \sum_{k} \mathbf{1}[\text{operator } o \text{ exits under } \Gamma'] = 0
      \quad \forall\, o,
      \label{eq:c4}
\end{align}
%
where $\Gamma' \equiv \{(i'_k, k)\}_{k \in \mathcal{K}}$ denotes a feasible
garage--operator assignment and $i'_k \in \mathcal{I}$ is the operator
assigned to garage $k$ under the counterfactual. The penalty term
$\hat{\theta}_1 \cdot \Delta\mathrm{YearsInc}_{k}(\Gamma')$ captures the
destruction of incumbent-specific investment that accompanies reallocation:
%
\begin{equation}
  \Delta\mathrm{YearsInc}_{k}(\Gamma') =
  \begin{cases}
    0 & \text{if } i'_k = i_k \quad \text{(incumbent retained)}, \\
    \mathrm{YearsInc}_{k} & \text{if } i'_k \neq i_k \quad \text{(operator
    changes)},
  \end{cases}
\end{equation}
%
and $i_k$ and $i'_k$ denote the operators assigned to garage $k$ under the
baseline $\Gamma$ and the counterfactual $\Gamma'$, respectively. No new
operators enter: $\mathcal{I}$ is held fixed, consistent with
constraint~\eqref{eq:c4}.
 
Constraint~\eqref{eq:c1} requires that aggregate proximity to own garages does
not deteriorate relative to baseline. Constraint~\eqref{eq:c4} rules out
operator exit, which would reduce competition and increase mark-ups.
Constraint~\eqref{eq:c3} merits further discussion. Since
$\sum_k \log(\mathrm{NrGarages})_{ik}$ is equivalent to
$\sum_i n_i \log n_i$, and $f(n) = n\log n$ is strictly convex, this sum is
minimised by equal market shares and increases sharply under concentration. The
constraint therefore permits equalising transfers but blocks concentrating
ones, providing a competition-policy safeguard analogous to the
Herfindahl--Hirschman Index.%
\footnote{To illustrate: with an initial distribution of $(2,4,8)$ garages,
$\sum_i n_i \log n_i \approx 23.6$. Equalising to $(4,5,5)$ lowers this to
$22.4$ (feasible), while concentrating to $(1,1,12)$ raises it to $29.8$
(infeasible under the constraint).}
The constraint is further motivated by the structural estimates: the negative
coefficient on $\log(\mathrm{NrGarages})$ implies that concentrating ownership
would mechanically raise measured welfare through that channel alone,
independently of any genuine efficiency gain; constraint~\eqref{eq:c3} rules
out this spurious source of improvement.
 
Because the objective depends on the full allocation through strategic network
interactions---via the coupled, non-separable terms $\mathrm{NrGarages}$,
$\mathrm{MinDurOwnGar}$, and $\mathrm{CompGar}$---the problem is a variant of
the quadratic assignment problem and NP-hard in general. The feasible space of
$8^{68} \approx 10^{61}$ candidate allocations rules out exhaustive search; we
therefore solve it by simulated annealing with multi-restart.%
\footnote{This approach is analogous to \citet{kreindler2024optimal}. The
algorithm exploits the fact that drive times between all garages and routes are
precomputed and fixed, so only the network variables of garages within a
ten-minute radius of each proposed reassignment need to be updated at each
step. Simulated annealing accepts welfare-improving swaps with probability one
and welfare-reducing swaps with a positive probability that declines according
to a cooling schedule, allowing escape from local optima early in the search.
We run three restarts, each initialised from the best state found so far. \florian{to update --- Florian I leave the fill in of the technical details to you, ok?}}
 
As in the urban-policy simulations of Section~\ref{SS:UP1}, we take the 2019
network as the benchmark, restricting attention to the \textcolor{red}{[XX]}
garages with observed capacity, positive bus assignments, and at least one
active route in 2019. We compute a \textbf{baseline} by assigning buses to
garages by proximity using equation~\eqref{eq:assignment_buses} given the
observed 2019 garage--operator allocation, and evaluate welfare
at~\eqref{eq:planner} under four increasingly rich specifications.
 
%% ------------------------------------------------------------------
\paragraph{v1 --- Operator network reorganisation.}
We iteratively propose swaps of garage ownership across operators and accept
welfare-improving reallocations subject to the constraints
in~\eqref{eq:planner}. Routes and the associated bus fleet follow the garage
to its new operator. Because bus-to-garage assignments are recomputed within
each operator's own garage network by proximity after each swap, dead-mile
covariates change; the garage-level variable $\mathrm{RouteStarts10Min}_k$
entering the structural utility is held fixed by construction, as it was
unobserved for most of the historical estimation sample. Comparing the
baseline to v1 therefore isolates the welfare effect of a different spatial
organisation of garage ownership: agglomeration benefits from denser operator
networks, reduced local monopoly rents from closer competitor proximity, and
the destruction of incumbent-specific investment captured by the
$\mathrm{YearsInc}$ term.
 
%% ------------------------------------------------------------------
\paragraph{v2 --- Dead-mile costs.}
We augment the objective with a garage--operator-specific dead-mile penalty
based on the full route-to-garage assignment $\mathcal{R}_k(\Gamma')$
available in the simulation. For each garage $k$ operated by operator $i$
under allocation $\Gamma'$, we subtract from $\hat{\Pi}_{ik}(\Gamma')$ the
penalty
%
\begin{equation}
  \frac{\beta}{\kappa} \cdot \frac{1}{|\mathcal{R}_k|}
  \sum_{r \in \mathcal{R}_k(\Gamma')} \mathrm{DeadMiles}_{ikr},
\end{equation}
%
where $\mathrm{DeadMiles}_{ikr}$ is the start--stop drive time (in minutes)
from garage $k$ to route $r$, $\beta = 0.019$ million pounds per minute of
empty driving time (from the dead-miles regressions), and $\kappa$ is the
utils-to-million-pounds conversion from Section~\ref{sec:counterfactuals}. To
rule out spurious welfare gains driven by routes being reassigned to more
distant garages---which mechanically inflates the penalty and hence raises
$W$---we impose the additional constraint
%
\begin{equation}
  \label{eq:c_dm}
  \sum_{k} \sum_{r \in \mathcal{R}_k(\Gamma')} \mathrm{DeadMiles}_{ikr}(\Gamma')
  \leq \lambda_{DM}
  \sum_{k} \sum_{r \in \mathcal{R}_k(\Gamma)} \mathrm{DeadMiles}_{ikr}(\Gamma).
\end{equation}
%
We report results for $\lambda_{DM} = 1$ (dead-mile costs weakly
non-increasing) and $\lambda_{DM} = 0.95$ (a five-percent reduction in total
dead miles is required). The key difference between v1 and v2 is that v2
captures the full operational cost of a dispersed garage network: operators
serving routes from geographically distant garages bear higher dead-mile costs
that are now explicitly penalised in the objective.
 
%% ------------------------------------------------------------------
\paragraph{v3 --- Bus fleet constraints.}
Versions v1 and v2 assume that operators can, in the medium to long run,
acquire sufficient buses to operate all routes that follow from a garage swap.
We now relax this assumption and consider what happens when an operator's bus
fleet is constrained at its 2019 baseline level. A proposed garage swap is
accepted only if the acquiring operator's baseline fleet is sufficient to cover
its existing route commitments plus the PVR of the routes inherited with the
garage. We consider two variants. In $W_3^{\mathrm{slow}}$, the fleet
constraint governs feasibility at entry but the operator's capacity expands
upon taking over a garage, allowing gradual fleet growth. In
$W_3^{\mathrm{fixed}}$, the constraint is strict: bus stocks never change and
operators always operate within their observed 2019 fleet. Both variants build
on the v2 objective with $\lambda_{DM} = 1$.
 
%% ------------------------------------------------------------------
\paragraph{v4 --- Internalising dead-mile externalities.}
Dead bus movements impose external costs on London residents through air
pollution, carbon emissions, noise, and congestion that are not borne by the
operating firm. Combining UK damage cost estimates for local pollution, noise,
congestion, and carbon, and assuming buses make on average two trips to and
from the route, we obtain a Pigouvian correction $\gamma$ equivalent to
approximately 20 percent of the private dead-mile penalty $\beta/\kappa$
already embedded in v2 and v3.%
\footnote{Using DEFRA central damage cost estimates for NO$_x$
({\pounds}10{,}193 per tonne) and PM$_{2.5}$ ({\pounds}114{,}411 per tonne),
the UK WebTAG central carbon value of {\pounds}248 per tonne CO$_2$, and
Transport for London unit cost estimates for noise ({\pounds}0.04 per km) and
urban congestion ({\pounds}0.05 per km), the combined social cost of empty bus
movements is approximately {\pounds}0.30 per km. Scaled to a five-year
contract period assuming two dead-mile trips per bus per day at an average
London bus speed of 18 km/h, and converted to utility units via $\kappa$, the
implied Pigouvian correction is
%
\begin{equation*}
  \gamma =
  \underbrace{\frac{c_{\mathrm{ext}}}{\kappa}}_{\text{utils/{\pounds}m}}
  \times
  \underbrace{n_{\mathrm{trips}} \cdot d_{\mathrm{contract}} \cdot s}_{%
    \text{km/(min$\cdot$bus$\cdot$contract)}}
\end{equation*}
%
where $c_{\mathrm{ext}} = {\pounds}0.299$ per km, $n_{\mathrm{trips}} = 2$,
$d_{\mathrm{contract}} = 1{,}825$ days, $s = 0.3$ km/min, and
$\kappa = 1.233$ {\pounds}m per util. This yields
$\gamma \approx 2.66 \times 10^{-4}$ utils per dead-mile minute per bus, or
$\gamma \times \bar{n}_{\mathrm{PVR}} \approx 0.20\cdot\theta_{DM}$ per
garage at mean PVR $\bar{n}_{\mathrm{PVR}} = 11.85$ buses per route.
\textcolor{red}{[UPDATE]}}
 
Version v4 augments the v3 objective with an explicit social benefit of
dead-mile reduction:
%
\begin{align}
  \label{eq:planner-v4}
  W_4 = \max_{\Gamma'} \; &
    -\sum_{i,k} \left[
      \hat{\Pi}_{ik}(\Gamma')
      - \frac{\beta}{\kappa}\,\frac{1}{|\mathcal{R}_k|}
        \sum_{r \in \mathcal{R}_k(\Gamma')} \mathrm{DeadMiles}_{ikr}
      + \hat{\theta}_1 \cdot \Delta\mathrm{YearsInc}_{k}(\Gamma')
    \right]
    + \gamma \cdot \bigl(\mathrm{DM}^0 - \mathrm{DM}(\Gamma')\bigr),
    \notag \\
  \text{subject to:}\quad
    & \eqref{eq:c1},\;\eqref{eq:c3},\;\eqref{eq:c4},\;\eqref{eq:c_dm},
\end{align}
%
where $\mathrm{DM}^0 = \sum_{k,r} \mathrm{PVR}_r \cdot d_{kr}$ and
$\mathrm{DM}(\Gamma') = \sum_{k}\sum_{r \in \mathcal{R}_k(\Gamma')}
\mathrm{PVR}_r \cdot d_{kr}$ denote total PVR-weighted dead miles at the
baseline and under the counterfactual, respectively. Unlike v2 and v3, where
the dead-mile penalty enters as a correction to measured operator profit---so
that a reduction in dead miles raises measured profits and \emph{reduces}
$W = -\sum_k \hat{\pi}_k$---version v4 adds a separate term that directly
rewards the planner for reducing total dead miles. The hard dead-mile
constraint from v2 and v3 is therefore dropped: the explicit social benefit
term already penalises dead-mile increases through the objective, making the
constraint redundant. Version v4 thus represents a pure Pigouvian
internalisation of the externality, trading off network reorganisation gains
against pollution abatement without an arbitrary constraint threshold.
 

Figure~\ref{fig:realloc-experiment} displays the main results;
Table~\ref{tab:welfare-realloc} decomposes the welfare effects by covariate.


\begin{figure}
    \caption{Results from garage reallocation experiment.\label{fig:realloc-experiment}} 
    \begin{adjustbox}{width= \textwidth}
        \includegraphics{\plots/map_realloc_panel.pdf}
    \end{adjustbox}
    \fnote{Each panel shows the resulting new allocation of garages if we maximize the social welfare function (i.e. minimize operator rents) under increasingly tight operational constraints. Panel a) imposes no constraints on the new allocation, panel b) required enough spare bus capacity, and panel c) adds an additional penalty for increased operational costs via additional dead miles from new garages. Details in appendix \ref{S:realloc-experiment}.}
\end{figure}

\begin{table}[!tb]
\caption{Decomposition of welfare effects from Garage Reallocation}\label{tab:welfare-realloc}
\centering
\resizebox{0.8\textwidth}{!}{%
\input{\tables/realloc_decomposition.tex}}
\fnote{Given the linearity of the structural model, we can compute for each covariate $k$ its change from baseline ($x_{kj}^0$) to optimal value as $\Delta x_{kj} = x_{kj}^* - x_{kj}^0$. The total change in welfare due to covariate $k$ is then $-\hat{\theta}_k \sum_j \Delta x_{kj}$, which is what the upper panel of the table shows.}
\end{table}

\section{Conclusions}\label{S:Conclusion}
This paper studies spatial rents and firm conduct in the London bus market, 
building on a new database of bus garages used by private firms for the 
operation of bus routes in the Greater London Area. The dataset contains 
details about the location and ownership of bus garages over time, starting 
from 1994 when the market was privatized along geographic lines, and is 
constructed by combining archival data from the London Omnibus Traction 
Society with information from industry hobbyist sources. The London bus 
procurement setting is particularly suitable for the analysis of spatial 
rents: with the demand side tightly regulated by TfL, the location of a 
garage in the network of routes and competing garages is the key determinant 
of a firm's profitability. We consider spatial rents broadly, capturing 
economies of density from operating routes from multiple nearby own garages, 
local monopoly rents from isolation relative to competitor garages, and dead 
mile transportation costs from proximity to the route network.

The key methodological contribution is a reversed discrete choice framework 
that exploits the exclusive use of garages to recast the firm location 
problem as one where garages choose operators, yielding a tractable 
multinomial estimation strategy that sidesteps the dimensionality and 
equilibrium multiplicity problems that arise in standard multi-player entry 
games. The garage value function is grounded in a structural procurement 
auction model that links spatial positioning to expected profitability, 
while remaining agnostic about competitive conduct at the auction 
stage---a flexibility that proves essential for the counterfactual analysis. 
A parsimonious model with low-dimensional representations of spatial 
centrality measures does remarkably well in approximating the data, as supported by model fit measures. The structural estimates 
reveal that both economies of density and local monopoly rents are 
economically meaningful determinants of garage value---an operator is 
indifferent between having its closest own garage three minutes nearer and 
having one fewer competitor garage within a ten-minute drive---and that 
spatial positioning, rather than historical patterns of use, explains the 
observed clustering of operators in certain parts of London.

Applying the estimates to counterfactual simulations, we find that private 
garage ownership generates an efficiency loss of \textsterling 34m-- 
\textsterling 51m per year---12 to 18 percent of total procurement 
costs---as incumbents retain garages that would be more valuable to 
competing operators. We further find that the pattern of unrealized 
profitable transactions correlates with market structure in ways consistent 
with tacit market-sharing: implied punishment values are highest during 
the stable oligopoly period of 2001--2008 when entry was absent, and 
correlate strongly with the spatial isolation of garages from competitor 
networks. These results suggest that the concerns voiced by the UK 
Competition Commission about potential market-sharing agreements in the 
London bus market are well-founded, though we stop short of claiming 
conclusive evidence of collusion. More broadly, the paper illustrates how 
the location choices of firms in regulated urban markets encode information 
about the spatial structure of competition that is not visible in price or 
quantity data alone. The framework applies to any market where firms 
maintain proprietary infrastructure, capacity is geographically fixed, and 
spatial proximity generates both cost synergies and competitive 
spillovers---conditions that characterize a wide range of industries from 
logistics and retail to utilities and public transport. A natural 
policy implication is that regulators in such markets should attend not 
only to pricing and procurement conduct but also to the spatial 
organization of firm infrastructure, which can generate and entrench 
market power in ways that standard regulatory tools are ill-equipped 
to address.


%% BIBLIOGRAPHY
\newpage
\bibliographystyle{aea}
\bibliography{../BusLocation_references} 




%% APPENDIX
\newpage
\clearpage
\appendix

\begin{center}{\large{[For Online Publication]}}\end{center}


\section{Additional results}

\setcounter{table}{0}
\renewcommand{\thetable}{A\arabic{table}}
\setcounter{figure}{0}
\renewcommand{\thefigure}{A\arabic{figure}}

\begin{table}[!h]
\caption{Examples of garages with central functions}\label{T:garage-syn}
\resizebox{\textwidth}{!}{%
\begin{tabular}{| l | l | L{\linewidth-4cm}|}
	\hline
	Operator&Garage 	&Description of synergy / task \\
	\hline
	Abellio	&Battersea Depot&	This large central London depot is home to some iconic London bus routes and is the centre for Abellio bus driver training and our iBus operations. Over 450 people are based here. Roles include operations, engineering, performance and HR.\\ \hline
	Abellio &	Walworth&	Welcome to our easily accessible, central depot in the heart of Camberwell. Over 500 people are based here, roles include operations, engineering and head office functions.\\ \hline
	Arriva London	&Enfield	&In recent times, Enfield has become a central part of Arriva London operations with the opening of the Accident Repair Centre, which also undertakes major refurbishments of the company’s fleet.\\ \hline
	Arriva London	&Norwood	&The garage previously supplied some buses and drivers for route 19 following the conversion from Routemaster operation in 2005.\\ \hline
	Arriva London	&Turnbridge Wells&	Lauren Edmonds from Arriva: ``It's fantastic that we're ready to move our engineering teams in at the new depot and we're expecting everything to be fully operational by the spring."\\ \hline
	Blue Triangle	&River Road&	Some buses are regularly shared with River Road garage. This garage also acts as a storage for the Here East bus fleet.\\ \hline
	East London	&West Ham	&The garage became fully operational in November 2009 taking over its own maintenance. The new garage is able to hold over 350 buses. It is the biggest bus garage in England and is the new location for Stagecoach London's head office and training centre.\\ \hline
	GoAhead	&Stockwell&	Buses are given a thorough check-over every 28 days or so by Stockwell's team of 35-odd engineers. Whenever you hear the announcement: ``This driver has been instructed to  hold the service", that call has been made manually, from a control room like the one at Stockwell. \\ \hline
	London General&	Merton&	Merton was also responsible for the maintenance of vehicles for route 200 in 1988/9 after the withdrawal of the Cityrama sightseeing company, whilst the route was operated from Sutton garage.\\ \hline
	London General	&Bexleyheath	&Even though Belvedere was a London General garage it operated as a satellite of London Central's Bexleyheath garage. All of Belvedere's engineering and fleet and staff management was done at or organised by Bexleyheath.\\ \hline
	RATP Dev London&	Edgware	&Edgeware is referenced as the company's ``operating base" and pictures show the workshop. \\ \hline
	RATP Dev London&	Harrow &	RATP Dev received planning approval for a three-bay double decker workshop at Harrow garage. The facility is to replace the existing maintenance area and allow for maintenance on the new fleet of electric buses.\\ \hline
	RATP Dev London	&Stamford Brook&	Flagship bus garage with RATP Dev Transit London headquarters. Location of comprehensive central engineering works.\\
	\hline
	\end{tabular}}
	\fnote{Garage task descriptions are quotes from the company's websites, collected during August 2024. Sources are available upon request.}
\end{table}
\clearpage

\begin{figure}[!h]
\centering
\includegraphics[width=0.7\textwidth]{\plots/ownership-count.pdf}
\caption{Number of ownership changes by garage}
\fnote{This map illustrates that different garages experience a different number of owners over time, in part depending on their location within the route network. Darker road segments imply a greater number or bus routes passing the same segment.\label{fig:ownership-count}}
\end{figure}

\begin{figure}[!h]
\centering
\includegraphics[width=0.8\textwidth]{\plots/ownership-timeseries.pdf}
\caption{Time series of Ownership Changes in garage data}
\fnote{We call \emph{Vacant} garages those which are vacant, abandoned or not yet built at a certain point in time. \emph{New} refers to a change of operator in a certain garage.\label{fig:ownership-TS}}
\end{figure}


\begin{figure}[!h]
\centering
\includegraphics[width=0.8\textwidth]{\plots/operator-garage-count.pdf}
\caption{Time series of \emph{Number of Garages} by operator}
\fnote{Number of garages owned by the eight largest operators by year, 
measured on the closest changedate prior to January 1st of each year. 
Garages owned by all remaining operators are aggregated as ``Other".} \label{fig:garages-TS-byoperator}
\end{figure}



\begin{figure}[!p]
\begin{minipage}{0.5\textwidth}
\centering
\subcaption{When owner changed}
\includegraphics[width = \textwidth]{\plots/nalt-by-year-mlogit.pdf}\label{F:alternativesa}
\end{minipage}%
\begin{minipage}{0.5\textwidth}
\centering
\subcaption{When owner stayed}
\includegraphics[width = \textwidth]{\plots/nalt-by-year-nonchoice.pdf}\label{F:alternativesb}
\end{minipage}%
\caption{Number of active operators that a garage can choose from}
\end{figure}


\begin{figure}[!h]
\centering
\includegraphics[width=0.7\textwidth]{\plots/GOF-entrants(6).pdf}
\caption{Model fit: ROC curves for the baseline estimates in Table \ref{T:strucestim}.}\label{F:GOFestimates}
\end{figure}

\begin{figure}
\centering
\includegraphics[width =0.7\textwidth]{\plots/choice-frequency/L-mlogit-model6.pdf}
\caption{Model fit: in-sample observed vs predicted entry frequency}
\fnote{Showing the in sample fit of the baseline estimates in Table \ref{T:strucestim}. Data choice frequency records the number of choice situations in data where an operator (including the Vacant operator) is the entrant, divided by the total number of choice situations (192), in all garage-change-dates where the operator changed. Model choice frequency records the share of these choice situations where the operator is assigned the highest predicted utility by the selected model.}\label{F:frequencies}
\end{figure}

\begin{figure}
\centering
\includegraphics[width = 0.8\textwidth]{\plots/choice-frequency/L-nonchoice-model6.pdf}
\caption{Out-of-sample observed vs model-predicted garage ownership for \emph{stayers}}
\fnote{Showing the out-of-sample predictions of the baseline estimates in Table \ref{T:strucestim} in comparison to observed garage ownership. Data choice frequency records the number of choice situations in data where an operator (including the Vacant operator) is the incumbent, divided by the total number of choice situations (66,964), in all garage-change-dates where the operator did not change. Model choice frequency records the share of these choice situations where the operator is assigned the highest predicted utility by the selected model.}
\label{F:frequenciesOOS}
\end{figure}

\begin{table}[!tb]
\caption{Structural estimates with alternative distance 
specifications\label{T:robustness_dist}}
\resizebox{\textwidth}{!}{%
\input{\tables/L-paper-robustness-distances.tex}}
\fnote{Column (1) is the baseline specification; variables are described 
in detail in Table~\ref{T:strucestim}. Columns (2)--(3) replace 
\textit{RouteStarts10Min} with alternative dead mile measures: the number 
of route stopover points within a 10-minute drive in column (2), and the 
number of stopover points within 500 metres in column (3), where a 
stopover point is defined as a bus stop where at least five routes 
intersect. Column (4) replaces \textit{MindurOwnGar} with 
\textit{AvgdurOwnGar}, the average drive time between garage $j$ and all 
other garages of the same operator. Column (5) replaces 
\textit{CompGar10Min} with \textit{MindurCompGar}, the minimum drive time 
between garage $j$ and the nearest competitor garage. Column (6) replaces 
\textit{CompGar10Min} with \textit{AvgdurCompGar}, the average drive time 
between garage $j$ and all competitor garages. Statistical significance: 
+ $p<0.1$, * $p<0.05$, ** $p<0.01$, *** $p<0.001$.}
\end{table}


\begin{table}[!tb]
\caption{Structural estimates with alternative incumbency and scale 
specifications\label{T:robustness_other}}
\resizebox{\textwidth}{!}{%
\input{\tables/L-paper-robustness-other.tex}}
\fnote{Column (1) is the baseline specification; variables are described 
in detail in Table~\ref{T:strucestim}. Column (2) replaces the linear incumbency term with 
$\exp(-\textit{YearsInc})$; column (3) omits the incumbency variable 
entirely. Column (4) replaces $\log(\textit{NrGarages})$ with 
\textit{NrGarages} in levels; column (5) omits the scale variable entirely. 
Column (6) omits the \textit{Nonvacant} intercept. Statistical significance: 
+ $p<0.1$, * $p<0.05$, ** $p<0.01$, *** $p<0.001$.}
\end{table}

\begin{table}
\caption{Private returns from capacity expansion, for garages with (near-)binding capacity}\label{tab:cap-expansion-binding}
\centering
\resizebox{0.83\textwidth}{!}{%
\input{\tables/capacity-expansion-ranking-binding.tex}}
\fnote{$\Delta \Pi$ (£m) is the change in private returns, in millions of pounds, from adding 50 buses to a garage relative to vacating the site (i.e.\ selling the plot at market value), holding the garage infrastructure and operator-bus network fixed. Private returns are defined in \eqref{eq:UP1} and decompose into two components: the \emph{Structural} term, which captures the average per-bus value of operating from that garage, and the \emph{DM correction}, which adjusts for the change in dead-mile costs as the 50 additional buses are drawn from marginal routes that may have different deadhead driving distances. $B^{\text{pre}}$ denotes the number of buses stationed at the garage before the capacity expansion; utilization is defined as $B^{\text{pre}}/\text{Capacity}$. Near-binding capacity defined as $\geq 80\%$ utilization at baseline. See Section~\ref{SS:urbanpolicy1} for further details.}
\end{table}


\begin{table}
\caption{Private returns from capacity expansion, for garages with excess capacity}\label{tab:cap-expansion-slack}
\centering
\resizebox{0.83\textwidth}{!}{%
\centering
\input{\tables/capacity-expansion-ranking-slack.tex}}
\fnote{$\Delta \Pi$ (£m) is the change in private returns, in millions of pounds, from adding 50 buses to a garage relative to vacating the site (i.e.\ selling the plot at market value), holding the garage infrastructure and operator-bus network fixed. Private returns are defined in \eqref{eq:UP1} and decompose into two components: the \emph{Structural} term, which captures the average per-bus value of operating from that garage, and the \emph{DM correction}, which adjusts for the change in dead-mile costs as the 50 additional buses are drawn from marginal routes that may have different deadhead driving distances. $B^{\text{pre}}$ denotes the number of buses stationed at the garage before the capacity expansion; utilization is defined as $B^{\text{pre}}/\text{Capacity}$. Excess capacity defined as $< 80\%$ utilization at baseline. See Section~\ref{SS:urbanpolicy1} for further details.}
\end{table}

\begin{table}[!tb]
\caption{Drivers of private returns to capacity expansion}\label{tab:cap-expansion-reg}
\centering
\resizebox{0.9\textwidth}{!}{%
\input{\tables/capacity-expansion-reg.tex}}
\fnote{OLS regressions with $\Delta \hat{\Pi}_{ij}$ \eqref{eq:UP1} as dependent variable, for all 62 garages in the sample with observed capacity that are not vacant and have buses stationed in 2019. Columns (1)--(3) use the nominal return to a 50-bus capacity expansion. Columns (4)--(6) deflate the dependent variable by the 2019 demeaned borough-level house price index, reducing the sample to the 55 garages located within the Greater London boundary for which borough-level house price data are available. Columns (7)--(9) adjust with more disaggregated ward-level house prices from HM Land Registry Price Paid Data from 2017, also demeaned across the observations in the sample.}
\end{table}


\clearpage
\section{Capacity orthogonality}\label{A:capacity}
\setcounter{table}{0}
\renewcommand{\thetable}{B\arabic{table}}
\setcounter{figure}{0}
\renewcommand{\thefigure}{B\arabic{figure}}
The structural garage value function in \eqref{eq:garagestruc} depends on 
garage capacity through the bus assignment rule \eqref{eq:assignment_buses}, 
but capacity is excluded from the reduced-form utility function 
\eqref{eq:garagevalue_detail} because it is garage-specific and therefore 
cancels from utility differences across operators in the reversed-choice 
framework. The identifying assumption is that capacity is orthogonal to the 
operator-specific network variables included in estimation. We test this 
assumption using the following procedure.

Capacity data (Peak Vehicle Requirement, PVR) is available for 61 out of 91 
garages in the estimation sample. For each garage, we compute the 
time-series average of each network variable across all changedates in the 
sample period, yielding one observation per garage. We then regress capacity 
on the averaged network variables and report simple correlations, partial 
correlations, and OLS regression results.

Table~\ref{T:capacity_correlation} reports the pairwise correlation matrix. 
The correlations between capacity and the network variables are small in 
magnitude: 0.088 with \textit{RouteStarts10Min}, $-0.116$ with 
\textit{MindurOwnGar}, 0.021 with \textit{CompGar10Min}, and $-0.200$ with 
\textit{NrGarages}. The correlation with the incumbency variable 
\textit{exp(YearsInc)} is negligible ($-0.002$). The largest off-diagonal 
correlations are between the network variables themselves rather than between 
any network variable and capacity, confirming that capacity is not the 
primary source of correlation in the data.

\begin{table}[!tb]
\caption{Pairwise correlation matrix: capacity and network 
variables\label{T:capacity_correlation}}
\resizebox{\textwidth}{!}{%
\input{\tables/capacity-correlation-matrix.tex}}
\fnote{Pairwise correlations computed across 61 garages with available 
PVR capacity data, using time-series averages of network variables over 
the full sample period 1995--2019. \textit{RouteStarts10Min}, 
\textit{MindurOwnGar}, \textit{CompGar10Min}, and \textit{NrGarages} 
are the network variables used in the baseline specification 
\eqref{eq:garagevalue_detail}. \textit{exp(YearsInc)} is the incumbency 
variable.}
\end{table}

Table~\ref{T:capacity_orthogonality} reports OLS regressions of capacity 
on the network variables, with and without pairwise interactions. In the 
main effects specification, the included regressors jointly explain less 
than 1 percent of the variation in capacity (adjusted $R^2 = 0.008$). 
Neither \textit{RouteStarts10Min}, \textit{MindurOwnGar}, nor 
\textit{CompGar10Min}---the variables most directly relevant to the spatial 
rent estimates---are statistically significant. The only significant 
association is with \textit{NrGarages} ($\hat{\beta} = -12.5$, $p<0.05$): 
larger operators tend to occupy smaller garages on average. Adding all pairwise interactions raises the 
adjusted $R^2$ to 0.083, driven by a significant 
\textit{CompGar10Min}$\times$\textit{MindurOwnGar} interaction 
($\hat{\beta} = -3.159$, $p<0.05$), suggesting that garages in competitively 
isolated but network-peripheral locations tend to be somewhat smaller. 
However, the low adjusted $R^2$ even in the interaction specification 
indicates this is a second-order pattern.

\begin{table}[!tb]
\caption{Orthogonality of garage capacity to network 
variables\label{T:capacity_orthogonality}}
\centering
\input{\tables/capacity-orthogonality.tex}
\fnote{OLS regressions of garage capacity (Peak Vehicle Requirement) on 
the network variables used in the structural estimation. Each variable is 
averaged over the full sample period for each garage. The main effects 
specification includes each network variable separately; the interactions 
specification adds all pairwise products. Standard errors in parentheses. 
Statistical significance: + $p<0.1$, * $p<0.05$, ** $p<0.01$, 
*** $p<0.001$.}
\end{table}

\begin{figure}[!tb]
\caption{Garage capacity versus other garage utility drivers\label{F:capacity_orthogonality}}
\includegraphics[width=\textwidth]{\plots/capacity-orthogonality.pdf}
\fnote{Each panel plots garage capacity (Peak Vehicle Requirement, 
vertical axis) against one network variable (horizontal axis), using 
garage-level time-series averages over 1995--2019 for 61 garages with 
available PVR data. The red line shows the OLS fit with a 95 percent 
confidence band. Panels correspond to \textit{RouteStarts10Min} (A), 
\textit{MindurOwnGar} (B), \textit{CompGar10Min} (C), and 
\textit{NrGarages} (D).}
\end{figure}

Figure~\ref{F:capacity_orthogonality} illustrates the regression results 
visually: the scatter plots show no systematic relationship between 
capacity and any of the four network variables, with flat OLS fits and 
wide confidence bands throughout. Taken together, the evidence supports 
the identifying assumption that excluding capacity from the discrete choice 
estimation does not bias the spatial rent estimates. 

\clearpage


\section{Can market-sharing explain unrealized garage 
transactions?}\label{A:collusion}
\setcounter{table}{0}
\renewcommand{\thetable}{C\arabic{table}}
\setcounter{figure}{0}
\renewcommand{\thefigure}{C\arabic{figure}}
This appendix explores whether unrealized profitable garage transactions 
show patterns consistent with collusive market-sharing. This line of inquiry 
is motivated by the UK Competition Commission's 2011 finding of operator 
conduct whereby firms avoid competing in each other's ``Core Territories,'' 
leading to geographic market segregation at an estimated annual cost of 
\textsterling 115m--\textsterling 305m to UK consumers and taxpayers. The 
investigation explicitly excluded London and Northern Ireland, but expressed 
concern that such behavior might be more widespread.

\subsection*{Setup}

For each garage-changedate pair where no ownership change took place, we 
ask whether another operator would have benefited more from the garage than 
the incumbent. Let $\mathcal{J}^{N}_t$ denote the set of garages where no 
ownership change took place during changedate $t$, and let $\mathcal{J}^N 
= \{\mathcal{J}_t^N\}_{t}$ collect these sets across all changedates. For 
each triple $(t, j \in \mathcal{J}^N, i \in \mathcal{N}_t)$, we compute 
the deterministic utility $\hat{f}_{ijt}$ based on the model estimates 
in \eqref{eq:garagevalue_detail}, excluding combinations where the outside 
option is either the choice alternative or the incumbent to avoid ascribing 
behavioral attributes to the market. The full set $\mathcal{J}^N$ contains 
66,964 operator-garage-changedate triples. Since $\Pi_{ijt}$ and 
$\Pi_{I_{jt}}$ are independently Gumbel-distributed with means 
$\hat{f}_{ijt}$ and $\hat{f}_{I_{jt}}$ and scale parameter 1, their 
difference follows a Logistic distribution with mean $\hat{f}_{ijt} - 
\hat{f}_{I_{jt}}$ and variance 1. The estimated utility excess is defined as
\begin{equation}\label{eq:psi_app}
\hat{\psi}_{ijt} = \big(\hat{f}_{ijt} - \hat{f}_{I_{jt}}\big) 
\mathbf{1}\big(\hat{f}_{ijt} - \hat{f}_{I_{jt}} \geq \tau\big)
\end{equation}
with $\tau = 0.721$ the 95th percentile of the standardized Logistic 
distribution. Whenever $\hat{\psi}_{ijt} > 0$, we can conclude with 95 
percent confidence that operator $i$ has higher utility for garage $j$ than 
the incumbent at that date, accounting for the unexplained component of the 
utility function. The main analysis focuses on the subset $\mathcal{J}^{Nmax}$, 
which retains only the operator with the highest predicted utility for each 
garage-changedate pair, yielding 6,308 triples and avoiding double-counting 
when multiple operators have positive $\hat{\psi}_{ijt}$ for the same 
garage-changedate pair.

Under a market-sharing interpretation, $\hat{\psi}_{ijt}$ provides a lower 
bound on the expected punishment required to deter the highest-value operator 
from entering: cases where a profitable transaction did not occur despite 
another operator having strictly higher utility are consistent with that 
operator being deterred by the threat of retaliation. We treat 
$\hat{\psi}:\hat{\psi}>0$ as measuring the height of the implied punishment and 
$\mathbf{1}(\hat{\psi}>0)$ as measuring its probability. We maintain that 
observed ownership changes reflect competitive behavior---so that the 
structural estimates are uncontaminated by collusion---and that absent a 
market-sharing agreement the highest-value firm enters, as established in 
Section~\ref{SS:entry}. We remain agnostic about the exact form of any 
agreement, which could range from reciprocal bilateral arrangements to 
unilateral punishment threats \citep[see][]{Porter1999, belleflamme2004market}.

\subsection*{Number of active firms}

Table~\ref{T:psi_reg_nalt} regresses $\hat{\psi}_{ijt}$ and 
$\mathbf{1}(\hat{\psi}_{ijt}>0)$ on the number of active operators at each 
changedate. The raw correlations require care in interpretation: more active 
operators generate more operator-garage-changedate triples, mechanically 
inflating both the count and magnitude of positive $\hat{\psi}$ values 
through an order statistic argument. The results are largely inconsistent 
with a simple market-sharing interpretation. For the height of punishment, 
the coefficient on the number of alternatives is positive and significant 
across all specifications including without year fixed effects (columns 1--4), 
implying that periods with more active firms are associated with larger 
implied punishment values---the opposite of what the theory predicts if 
more firms make collusion harder to sustain. For the probability of a 
positive $\hat{\psi}$, the coefficient is negative without year fixed 
effects in the full sample (column 5, $p<0.001$), consistent with 
coordination becoming harder as more firms enter, but turns positive and 
significant once year fixed effects are included or when restricting to 
$\mathcal{J}^{Nmax}$ (columns 6--8). With a model specification that we circulated in a working paper, we found a negative and significant coefficient on the number of 
alternatives for the probability of punishment without year fixed effects, 
and an insignificant coefficient for the height---a pattern more consistent 
with the market-sharing prediction, though still not robust across all 
samples. Results for this alternative specification are available 
from the authors upon request. These sign reversals across specifications 
and the general sensitivity of the results to the structural model 
parameterization caution against placing weight on the number-of-firms 
results as evidence of market-sharing behavior. 

\begin{table}[!tb]
\centering
\caption{Regressions of punishment on the number of 
alternatives\label{T:psi_reg_nalt}}
\begin{adjustbox}{max width = \textwidth}
\input{\tables/psi-reg-nalt.tex}
\end{adjustbox}
\fnote{Statistical significance: + $p<0.1$, * $p<0.05$, ** $p<0.01$, 
*** $p<0.001$. $\mathcal{J}^{N}$ denotes the full sample of 66,964 
operator-garage-changedate triples where no ownership change took place; 
$\mathcal{J}^{Nmax}$ retains only the operator with the highest predicted 
utility for each garage-changedate pair, yielding 6,308 triples. 
$\hat{\psi}:\hat{\psi}>0$ denotes regressions explaining the magnitude of the implied 
punishment conditional on it being positive; $\mathbf{1}(\hat{\psi}>0)$ denotes 
regressions explaining the probability of a positive implied punishment.}
\end{table}


Figure~\ref{F:TimeSeries} plots yearly averages of $\hat{\psi}$ and 
$\mathbf{1}(\hat{\psi}>0)$ alongside years of significant entry, defined 
as the entry of at least one operator who remained active for at least five 
years. In the full sample (panels a--b), the probability of punishment 
broadly tracks the entry pattern---declining during high-entry periods and 
rising during the stable oligopoly period of 2001--2008---while the height 
of punishment also peaks in the low-entry period but with less pronounced 
variation. In the restricted sample $\mathcal{J}^{Nmax}$ (panels c--d), 
the height of punishment is more clearly elevated during the stable period, 
but the punishment probability declines steadily over time and does not 
increase during the low-entry period, contrary to the market-sharing 
prediction. Overall, the time-series patterns are mixed and do not 
uniformly support the collusion interpretation.


\begin{figure}[!tb]
\caption{Punishment and relation to firm entry\label{F:TimeSeries}}
\centering
\begin{minipage}{0.5\textwidth}
\centering
\subcaption{mean $\hat{\psi}(\hat{\psi}>0)$, all $ijt$}
\includegraphics[width=\textwidth]{\plots/time-series-psi-entry/mean-yearly-selected-all.pdf}
\end{minipage}%
\begin{minipage}{0.5\textwidth}
\centering
\subcaption{$1(\hat{\psi}>0)$, all $ijt$}
\includegraphics[width=\textwidth]{\plots/time-series-psi-entry/nonzero-yearly-selected-all.pdf}
\end{minipage}
\begin{minipage}{0.5\textwidth}
\centering
\subcaption{mean $\hat{\psi}(\hat{\psi}>0)$, max-$i$ per $jt$}
\includegraphics[width=\textwidth]{\plots/time-series-psi-entry/mean-yearly-selected-highest.pdf}
\end{minipage}%
\begin{minipage}{0.5\textwidth}
\centering
\subcaption{$1(\hat{\psi}>0)$, max-$i$ per $jt$}
\includegraphics[width=\textwidth]{\plots/time-series-psi-entry/nonzero-yearly-selected-highest.pdf}
\end{minipage}
\fnote{Yearly averages of the estimated punishment when positive (panels a 
and c) and the punishment probability (panels b and d). Panels (a) and (b) 
average across all entries of $\mathcal{J}^N$; panels (c) and (d) restrict 
to the highest-utility operator for each garage-changedate pair 
($\mathcal{J}^{Nmax}$). Vertical bars denote years with significant entry, 
defined as the entry of at least one operator who remained active for at 
least five years.}
\end{figure}

\subsection*{Spatial isolation from competitor garages}

Table~\ref{T:psi_reg_dist2comp} relates $\hat{\psi}_{ijt}$ to the spatial 
isolation of garages from competitor networks. Garages more isolated from 
competing garages---measured by either the minimum or average drive time to 
a competitor garage---are associated with both a higher probability of 
positive $\hat{\psi}$ and a larger magnitude conditional on being positive, 
across all specifications and robust to year fixed effects. The competitor 
distance variables explain 6--9 percent of the variation in $\hat{\psi}$. 
This pattern is directly consistent with the local monopoly rent mechanism 
estimated in Section~\ref{S:Results}: operators with more isolated garages 
earn higher expected route auction revenues and therefore face stronger 
incentives to protect their position from entry. The fact that $\hat{\psi}$ 
correlates strongly with the spatial rent variables that drive the structural 
estimates supports the interpretation that the unrealized transactions reflect 
economically meaningful forces rather than estimation noise, and this finding 
is robust to the structural model specification.

\begin{table}[!tb]
\centering
\caption{Regressions of punishment on distance to competitors' 
garages\label{T:psi_reg_dist2comp}}
\begin{adjustbox}{max width = \textwidth}
\input{\tables/psi-reg-dist2comp.tex}
\end{adjustbox}
\fnote{Statistical significance: + $p<0.1$, * $p<0.05$, ** $p<0.01$, 
*** $p<0.001$. $\mathcal{J}^{N}$ denotes the full sample of 66,964 
operator-garage-changedate triples where no ownership change took place; 
$\mathcal{J}^{Nmax}$ retains only the operator with the highest predicted 
utility for each garage-changedate pair, yielding 6,308 triples. 
$\hat{\psi}:\hat{\psi}>0$ denotes regressions explaining the magnitude of the implied 
punishment conditional on it being positive; $\mathbf{1}(\hat{\psi}>0)$ denotes 
regressions explaining the probability of a positive implied punishment.}
\end{table}

Figure~\ref{F:frequenciesOOS} offers a suggestive picture of the winners 
and losers from a reallocation of garages to their highest-utility operators. 
Smaller operators retain a similarly modest footprint in the counterfactual, 
confirming that the results are not driven by fringe operators becoming 
dominant players. At the same time, virtually all smaller operators expand 
somewhat. The most striking reallocation concerns Arriva, which shrinks from 
the largest operator---accounting for roughly 25 percent of entries in the 
observed network---to a fringe player in the counterfactual. Conversely, 
GoAhead, ComfortDelGro, and Abellio all more than double their garage shares. 
To the extent that these reassignments reflect garages blocked by 
market-sharing agreements, this pattern is consistent with Arriva being a 
net beneficiary of the collusive arrangement while GoAhead, ComfortDelGro, 
and Abellio are the operators most constrained by it.

The spatial rent model is identified entirely from observed ownership 
transitions, without using procurement bid data, and is agnostic about 
competitive conduct at the auction stage. The only assumption required is 
that transactions that do take place allocate garages to the operator that 
values them most---the model makes no use of information about transactions 
that did not occur. Operator-specific deviations from the model's predictions 
likely reflect a combination of unobserved operational considerations and 
firm-specific strategic behavior that the five network variables do not fully 
capture, and we stop short of claiming these results as unequivocal evidence 
of collusion.

\clearpage
\section{Validating the route-garage assignment equation}
\label{A:assignment_validation}
\setcounter{table}{0}
\renewcommand{\thetable}{D\arabic{table}}
\setcounter{figure}{0}
\renewcommand{\thefigure}{D\arabic{figure}}

This appendix validates the bus-garage assignment equation~\eqref{eq:assignment_buses} against observed
route-garage assignments in the last tender round. The sample consists of
all routes for which an operating garage is recorded in the tender data and
for which the observed garage has a known PVR capacity, yielding 584
route-operator pairs across eight operators.
 
For each operator $i$, let $\mathcal{J}_i$ denote the set of garages owned
by $i$ with known PVR capacity $K_j$, and let $\mathcal{R}_i$ denote the
corresponding set of routes with known vehicle requirement $q_r$. All pairs
$(j,r)\in\mathcal{J}_i\times\mathcal{R}_i$ are ordered by drive time from
garage $j$ to route $r$'s start stop ascending, with ties broken by $q_r$
descending. The algorithm iterates through this queue: route $r$ is assigned
to garage $j$ if $r$ is not yet assigned and $j$'s remaining capacity
$\tilde{K}_j\geq q_r$, after which $\tilde{K}_j$ is decremented by $q_r$.
Each route is assigned exclusively to a single garage; routes for which no
garage owned by operator $i$ retains sufficient remaining capacity are left
unassigned. The procedure runs separately for each operator, so that garages
are never shared across operators.
 
Two fit criteria are reported. \textit{Exact match} requires that the
predicted and observed garage codes agree. The \textit{within 2 minutes}
criterion relaxes this to a cost-equivalence standard: the prediction is
deemed acceptable if the predicted garage's dead-mileage to the route
differs from the observed garage's by at most two minutes. This tolerance
is evaluated separately under the three dead-mile measures used throughout
the paper---\emph{Dead Miles (Start)}, \emph{Dead Miles (StartStop)}, and
\emph{Dead Miles (Stopover)}---so that a mismatch in garage identity that
carries no meaningful cost implication is not counted as a failure.
 
Tables~\ref{tab:assignment_validation} and~\ref{tab:assignment_validation_30perc}
report results by operator under two capacity assumptions. Each table contains
two panels. The upper panel (\textit{Exclusive route--garage mapping}) imposes
that each route's entire vehicle requirement is met from a single garage,
consistent with how TfL routes are contracted as indivisible units. The lower
panel (\textit{Routes can spread over garages}) relaxes this, allowing a
route's buses to be distributed across multiple garages of the same operator,
with exact match defined as the observed garage appearing anywhere in the
predicted allocation. Table~\ref{tab:assignment_validation} uses observed
garage PVR capacities; Table~\ref{tab:assignment_validation_30perc} inflates
all capacities by 30 percent.\footnote{This is done as we observe that the baseline assignment of routes to garages implies that, for the garages in the sample, the total route PVR exceeds the total garage capacity PVR by 30 percent, so we want to understand whether this has consequences for the fit. Total route PVR can exceed garage capacity because of the way routes are run---e.g., night routes and day routes alternating, and buses rotating through the garage, or garages sometimes being parked too the brim--- or measurement error.}
 
The baseline specification---exclusive route assignment at observed
capacities---yields an overall exact match rate of 47.8 percent, rising to
63.4 percent (\emph{Dead Miles (Start)}), 71.7 percent
(\emph{Dead Miles (StartStop)}), and 64.4 percent (\emph{Dead Miles
(Stopover)}) under the within-two-minute criterion. Allowing route splitting
does not improve fit and slightly reduces exact match (46.7 percent overall),
consistent with capacity fragmentation: greedily placing partial bus
requirements across garages leaves small residual capacities that cannot
accommodate subsequent routes. Expanding capacity by 30 percent raises exact
match modestly to 49.1 percent (exclusive) and 48.6 percent (splitting),
confirming that binding capacity constraints account for a share of
misassignments but are not the primary source of prediction error.
Operator-level rates range from 88.9 percent exact matches for TowerTransit
to 26--28 percent for Stagecoach, whose route PVR requirements substantially
exceed its recorded garage capacity. The within-two-minute rates are
consistently higher than exact match rates across all operators and
specifications, indicating that prediction errors are predominantly small in
dead-mileage terms. We adopt the baseline exclusive specification for all
counterfactual exercises, as it is institutionally appropriate and performs no
worse than the alternatives.
 
\begin{table}[!tb]
\caption{Validation of the route-garage assignment equation}\label{tab:assignment_validation}
\centering
\input{\tables/assignment-validation-baseline.tex}
\fnote{The assignment equation~\eqref{eq:assignment_buses} assigns routes
exclusively to garages owned by the same operator, sequentially by minimum
drive time to the route start stop subject to garage PVR capacity. Routes
are treated as indivisible: the full vehicle requirement must be met from a
single garage. The sample covers the last tender round, restricted to routes
with an observed operating garage of known capacity (584 routes, 8
operators). \textit{Exclusive route--garage mapping}: each route is assigned
to at most one garage; exact match requires agreement of predicted and
observed garage codes. \textit{Routes can spread over garages}: a route's
buses may be distributed across multiple garages of the same operator; exact
match requires the observed garage to appear in the predicted allocation.
\textit{Within 2 min} = the predicted garage's dead-mileage differs from the
observed garage's by at most 2 minutes, evaluated under three measures: drive
time to the route start stop (\emph{Dead Miles (Start)}), average of start
and end stop (\emph{Dead Miles (StartStop)}), and mean over intermediate
stopover stops (\emph{Dead Miles (Stopover)}).}
\end{table}
 
\begin{table}[!tb]
\caption{Validation of the route-garage assignment equation (garage capacity $+30\%$)}\label{tab:assignment_validation_30perc}
\centering
\input{\tables/assignment-validation-capacity-x1.3.tex}
\fnote{As Table~\ref{tab:assignment_validation}, except that all garage PVR
capacities are scaled up by 30 percent prior to assignment. This robustness
check assesses sensitivity to capacity measurement error and to the share of
unassigned routes attributable to binding capacity constraints rather than
to geographic mismatch.}
\end{table}


\clearpage


\section{Private returns to capacity expansion---simulation details}\label{A:Capacity_expansion_details}
This section provides additiconal details underlying the measurement of the private returns to investing in capacity (Section \ref{SS:UP1}). The exercise computes for each garage
$j$ owned by operator $i$, the marginal route(s) drawn in by increasing the capacity of that garage by 50 buses, and combines this with the estimated garage valuations
$\hat{\Pi}_{ij}$ to recover the per-bus value of additional capacity. It relies on 
the assignment technology of equation~\eqref{eq:assignment_buses} to 
translate a capacity change into a predicted change in route stationing. 
The empirical validity of this step is supported by the validation in 
Appendix~\ref{A:assignment_validation}, where the same equation recovers 
the correct operating garage for 46 percent of routes exactly, and for 
62--70 percent of routes once a two-minute dead-mileage tolerance is 
applied---within which cost differences between predicted and observed 
assignments are negligible.

Measuring returns to investment in monetary terms requires translating 
structural utility into procurement costs. Since garage transaction prices 
are not observed, we match total predicted utils to observed TfL auction 
revenues in the same period, yielding an implied revenue per util $\kappa$ 
from the logsum formula applied to the preferred specification. Since the 
garage value function \eqref{eq:garagestruc} is derived from expected 
auction profits, utility ratios map to profit ratios up to the logit scaling 
factor $\sigma$, which is common across operators and garages under 
homoskedasticity. Using average annual TfL route auction revenues 
$\bar{R} = \text{\textsterling}273\text{m}$ over 2004--2019:
$$\kappa = \bar{R} \Big/ \sum_{ijt} \exp(\hat{f}_{ijt}) = 
\text{\textsterling}1{,}232{,}657 \text{ per util}$$
where the denominator is the logsum of predicted utilities across all 
operator-garage-changedate triples in the estimation sample during the same period.\footnote{In 
discrete choice models with an identified price coefficient $\alpha$, utils 
convert to monetary values at rate $1/\alpha$. Here $\alpha$ is not 
identified from garage choices alone, so we exploit the structural link 
between garage utility and auction profits: utility ratios map to profit 
ratios with $\sigma$ cancelling in the ratio.} We apply $\kappa$ throughout 
the counterfactual analysis to ensure structural utility differences and 
reduced-form dead miles costs are expressed on a common monetary scale 
anchored to observed TfL procurement payments.

\begin{figure}[!tb]
\caption{Role of capacity in Urban policy 1}\label{F:pi_bus_capacity}
\includegraphics[width=\textwidth]{\plots/pi-bus-capacity.pdf}
\end{figure}

For the structural term and dead miles correction in \eqref{eq:UP1} to be 
internally consistent, both must measure the change in operator profit per 
additional bus in the same monetary units. The dead miles correction is in 
\textsterling m by construction from procurement bid data; the structural 
term is converted via $\kappa$ derived from the same revenues, so both 
share a common monetary scale under the assumption that TfL auction revenues 
are the sole source of operator value. The remaining identification issue 
concerns the normalization of $\hat{\pi}_{ij}^{bus}$: since structural utility 
is measured relative to the outside option, a negative structural term means 
only that the operator values the garage less than the market, not that 
capacity expansion destroys value in absolute terms. The sum 
$\Delta\hat{\Pi}_{ij}$ therefore identifies a valid ranking across garages 
as long as the outside option value per bus does not vary systematically 
with the network variables that drive $\hat{\pi}_{ij}^{bus}$ --- precisely the 
orthogonality condition that underlies the structural estimates and is 
verified empirically in Appendix~\ref{A:capacity}.\footnote{We also assess the two components separately, which does not require them to lie on the same scale.}

Moreover, it must hold that 
$\hat{\pi}_{ij}^{bus}$ reflects the average per-bus profit at 
the baseline dead miles, so that the correction term adjusts for deviations 
of marginal routes from this baseline. A necessary condition is that 
$\hat{\pi}_{ij}^{bus}$ is uncorrelated with garage capacity after conditioning 
on network variables — otherwise the division by $B_{jt}$ introduces a 
mechanical bias. We verify this empirically (see Figure~\ref{F:pi_bus_capacity}). The per-bus value  negatively correlated with capacity in the full sample, but capacity explains less than five percent of the variation. Moreover, when excluding three outliers (small garages with much higher $\hat{\pi}_{ij}^{bus}$ than the other garages), capacity is positively correlated with per-bus value. This confirms that the per-bus 
measure reflects genuine network position heterogeneity rather than a 
capacity artifact, so that we can rely on $\hat{\pi}_{ij}^{bus}$ without mechanically disadvantaging larger garages.




\section{Constrained Welfare Reallocation Experiments\label{S:realloc-experiment}}

\subsection{Setup}

Let $\mathcal{J}$ denote the set of bus garages and $\mathcal{I}$ the set of
operators active in the 2019 market snapshot.  An \emph{allocation}
$\alpha: \mathcal{J} \to \mathcal{I}$ assigns each garage to exactly one
operator. We refer to the structural model defined in \eqref{eq:garagevalue_detail}, and we will recomputed operator-dependent covariates (subscript $i$) for
each proposed allocation $\alpha$.  Structural parameters
$(\theta_1,\ldots,\theta_6)$ are held fixed throughout all counterfactual experiments.

The social planner's objective is to minimise total industry profit:
%
\begin{equation}
  \min_{\alpha:\,\mathcal{J}\to\mathcal{I}} \;
  \sum_{j \in \mathcal{J}} \tilde{\Pi}_{\alpha(j),\,j}(\alpha).
  \label{eq:planner}
\end{equation}

which is the negative of our welfare function in the optimization algorithm.

\begin{equation}
    \max_{\alpha:\,\mathcal{J}\to\mathcal{I}} \; W(\alpha) \equiv -\sum_{j \in \mathcal{J}} \tilde{\Pi}_{\alpha(j),\,j}(\alpha)\label{eq:planner-welfare}
\end{equation}

%
We solve \eqref{eq:planner-welfare} by simulated annealing (SA).  Three
experiment variants impose progressively tighter operational constraints.

\subsection{Layer~0 - Unconstrained Reallocation}

The reference experiment imposes no operational constraints.  A proposed
swap of garage $j$ from incumbent $I$ to operator $i$ is always
admissible.  Buses and route contracts are treated as transferring
costlessly with the garage, so no adjustment to \eqref{eq:pi_hat_alloc}
is required.  This variant corresponds to a regulator that can costlessly
reassign full bundles (infrastructure, fleet, and route licences) to
whoever values them most.

\subsection{Layer~1 - Fixed Bus Stock}

Layer~1 relaxes the frictionless-transfer assumption by fixing each
operator's total bus fleet at its baseline level.

\paragraph{Bus stock.}
Each operator $i$'s fleet is set to
$\bar{B}_i = \sum_{j \in \mathcal{J}^0_i} \text{cap}(j)$,
where $\mathcal{J}^0_i$ is operator $i$'s baseline garage set and
$\text{cap}(j)$ is the PVR capacity of garage $j$.  Capacity is measured
as the maximum of the physical garage capacity (from survey data) and the
total PVR of routes assigned to $j$ in the last tender round.
$\bar{B}_i$ is fixed for the duration of the experiment.

\paragraph{Route inheritance.}
When garage $j$ transfers from incumbent $I$ to operator $i$, all routes
currently served from $j$ are reassigned to $i$.  Formally, if
$\mathcal{R}_j(\alpha)$ denotes the set of routes assigned to garage $j$
under allocation $\alpha$ (determined by a greedy proximity rule,
equation \eqref{eq:assignment_buses} in the main text), then after the
transfer:
\[
  \alpha(r) \leftarrow i \quad \forall\, r \in \mathcal{R}_j(\alpha).
\]
Route-to-garage assignments for $i$ and $I$ are then recomputed by
proximity, holding all other operators' assignments fixed.

\paragraph{Feasibility constraint.}
A proposed swap $j: I \to i$ is admissible only if operator $i$'s spare
bus capacity covers the inherited PVR:
\begin{equation}
  \bar{B}_i \;\geq\; \text{PVR}_i(\alpha) + \sum_{r \in \mathcal{R}_j(\alpha)} \text{pvr}(r),
  \label{eq:feasibility_L1}
\end{equation}
where $\text{PVR}_i(\alpha) = \sum_{r:\,\alpha(r)=i} \text{pvr}(r)$ is
operator $i$'s current total PVR commitment and $\text{pvr}(r)$ is the
peak vehicle requirement of route $r$.  Proposals violating
\eqref{eq:feasibility_L1} are rejected.

\subsection{Layer~2 - Fixed Bus Stock with Dead-Mile Penalty}

Layer~2 retains the feasibility constraint~\eqref{eq:feasibility_L1} and
adds an explicit operational cost for routes inherited at a garage whose
endpoint geography is misaligned with operator $i$'s existing network.

\paragraph{Penalty term.}
Let $\text{dm}(j, r)$ denote the OSRM drive time (minutes) from garage
$j$ to the start--stop endpoints of route $r$.  The structural utility at
garage $j$ for the receiving operator $i$ is adjusted to
\begin{equation}
  \tilde{\Pi}^{(2)}_{ij}(\alpha)
    = \tilde{\Pi}_{ij}(\alpha)
    + \theta_{\text{DM}} \cdot
        \frac{1}{|\mathcal{R}_j(\alpha)|}
        \sum_{r \in \mathcal{R}_j(\alpha)} \text{dm}(j, r),
  \label{eq:pi_hat_L2}
\end{equation}
where $\theta_{\text{DM}} \leq 0$.  The term represents the per-bus-unit
cost of operating routes that come with garage $j$ but whose endpoints
are distant, averaged across those routes.

\paragraph{Calibration of $\theta_{\text{DM}}$.}
The monetary bridge (equation~\ref{eq:monetize} in the main text) links
the logit utility scale to pounds:
\[
  \pi^{\text{mon}}_m = \hat{\pi}^* \cdot |\zeta| \cdot B,
\]
where $\zeta = \partial\,\text{AcceptedBid}/\partial\,\text{DeadMiles}$
is the dead-miles coefficient from the preferred dead-miles regression
(equation~\ref{eq:dead_miles_reg}, model~(2): $\zeta \approx -\pounds19{,}000$
per dead-mile-minute), and $B$ is garage capacity in buses.  Converting
the reduced-form estimate to structural utility units:
\begin{equation} 
  \theta_{\text{DM}} = \frac{\zeta}{|\hat{\zeta}| \cdot \bar{B}},
  \label{eq:theta_DM}
\end{equation}
where $\bar{B}$ is a representative garage capacity (median across the
sample).  Equation~\eqref{eq:theta_DM} ensures that one additional
dead-mile-minute translates into the same reduction in $\tilde{\Pi}^{(2)}$
as one pound of bid cost implies in the structural model.

\paragraph{Locality.}
The penalty in \eqref{eq:pi_hat_L2} is specific to garage $j$ and the
routes $\mathcal{R}_j(\alpha)$ that follow it.  It does not affect the
profit terms of any other garage, so the SA local-update structure (where
a single swap modifies only a small subset of the objective) is
preserved.

\subsection{Computational Details}

All three variants are solved by Metropolis--Hastings simulated annealing
with geometric cooling:
\[
  T_k = T_0 \cdot \alpha^k, \qquad \alpha = 0.9999,
\]
where $T_0$ is calibrated so that approximately 50\% of feasible proposals
are accepted at the start of the search.  For Layer~1 and Layer~2,
$T_0$ is estimated from the empirical distribution of $|\Delta W|$ over
feasible probe swaps.

At each iteration a garage is drawn uniformly at random and reassigned to
a uniformly random alternative operator.  A proposal is accepted with
probability $\min(1, \exp(\Delta W / T_k))$, subject to the feasibility
gate~\eqref{eq:feasibility_L1} for Layers~1 and~2.  The best state
(highest $W$) encountered across all iterations is returned.  Results are
verified with multi-restart runs; the best solution across restarts is
reported.

Table~\ref{tab:experiment_variants} summarises the three variants.

\begin{table}[h]
\centering
\caption{Experiment variants: assumptions on operational constraints}
\label{tab:experiment_variants}
\begin{tabular}{lccc}
\hline
 & Layer~0 & Layer~1 & Layer~2 \\
\hline
Buses transfer with garage       & Yes & No  & No  \\
Routes follow garage (Reading~B) & ---  & Yes & Yes \\
Bus-stock feasibility gate       & No  & Yes & Yes \\
Dead-mile penalty $\theta_{\text{DM}}$ & No & No & Yes \\
\hline
\end{tabular}
\end{table}

\noindent
Layer~0 provides an upper bound on achievable welfare gains: it asks what
the optimal allocation would be if the regulator could reassign full
operational bundles at zero cost.  Layer~1 introduces the realistic
constraint that operators cannot instantly field buses they do not own.
Layer~2 prices the operational friction of inheriting routes whose
geography does not fit the new operator's network.

\subsection{Welfare Decomposition by Covariate Channel}

Because the structural utility \eqref{eq:pi_hat_alloc} is linear in
covariates, the total welfare gain $\Delta W = W^* - W^0$ can be
decomposed exactly into additive contributions from each covariate
channel.

\paragraph{Setup.}
Let $x_{kj}^0$ denote the value of covariate $k$ at garage $j$ under the
2019 baseline allocation, and $x_{kj}^*$ its value under the SA-optimal
allocation.  Define the per-garage change
\[
  \Delta x_{kj} = x_{kj}^* - x_{kj}^0.
\]
Note that operator-level covariates ($\text{MindurOwnGar}$,
$\text{NrGarages}$, $\text{CompGar10Min}$) are recomputed for each
allocation, so $\Delta x_{kj}$ captures both direct reassignment effects
and the network-wide externalities that a swap induces on all garages
belonging to the operators involved.

\paragraph{Decomposition formula.}
The planner's objective is $W = -\sum_j \tilde{\Pi}_{\alpha(j),j}(\alpha)$,
so
\[
  \Delta W
  = W^* - W^0
  = -\sum_{j} \Delta\tilde{\Pi}_j
  = -\sum_{j} \sum_{k} \hat{\theta}_k\, \Delta x_{kj}
  = \sum_{k} \underbrace{\Bigl(-\hat{\theta}_k \sum_{j} \Delta x_{kj}\Bigr)}_{\text{channel-}k\text{ contribution}}.
\]
Each term $-\hat{\theta}_k \sum_j \Delta x_{kj}$ is the welfare gain
attributable to covariate channel $k$, and these terms sum exactly to
$\Delta W$.  A positive entry means the optimal allocation moves
covariate $k$ in the direction that reduces total operator profit: for
channels with $\hat{\theta}_k < 0$ (e.g.\ own proximity, competition)
this requires $\sum_j \Delta x_{kj} > 0$; for channels with
$\hat{\theta}_k > 0$ (e.g.\ incumbency decay) it requires
$\sum_j \Delta x_{kj} < 0$.

\paragraph{HHI.}
The Herfindahl--Hirschman index is computed on garage counts:
\[
  \text{HHI}(\alpha) = \sum_{i \in \mathcal{I}}
    \left(\frac{|\{j : \alpha(j) = i\}|}{|\mathcal{J}|}\right)^2.
\]
It is reported for the baseline and optimal allocations to summarise how
market concentration changes under each experiment layer; it is not an
input to the planner's objective.

\begin{figure}[!tb]
\caption{Solver paths in Reallocation experiment}\label{F:solver-paths}
\includegraphics[width=\textwidth]{\plots/convergence_realloc_layers.pdf}
\end{figure}

\end{document}